\documentclass[12pt,letterpaper]{article}
\usepackage[utf8]{inputenc}
\usepackage[spanish]{babel}
\usepackage{amsmath}
\usepackage{amsfonts}
\usepackage{amssymb}
\usepackage{tikz}
\usetikzlibrary{babel, arrows.meta} %
\usepackage{pgfplots}
\usepackage{listings}
\usepackage{xcolor}

\pgfplotsset{compat=1.18}

\lstset{
    language=Python,
    basicstyle=\ttfamily\small,
    keywordstyle=\color{blue}\bfseries,
    stringstyle=\color{red},
    showstringspaces=false,
    numbers=left,
    numberstyle=\tiny\color{gray},
    breaklines=true,
    frame=single,
    backgroundcolor=\color{white},
    rulecolor=\color{black}
}

\begin{document}

\subsection*{Enunciado}
¿Cuál es el castigo del fraude científico en la comunidad científica?

\subsection*{Solución}

\subsubsection*{Fundamentos de la Ética Científica}
En el estudio de las ciencias, y tal como lo expone Paul G. Hewitt en su obra "Física Conceptual", la ciencia no solo es un conjunto de conocimientos, sino también una actividad humana basada en la confianza y la honestidad. Los científicos confían en los datos y descubrimientos reportados por sus colegas para avanzar en la investigación. A diferencia de otros campos donde el engaño podría pasar desapercibido o tener penalizaciones temporales, la ciencia tiene un mecanismo de autocorrección muy riguroso.

\subsubsection*{Análisis del Fraude Científico}
El fraude científico, que incluye la falsificación de datos, la invención de resultados o el plagio, atenta directamente contra el pilar fundamental del método científico: la reproducibilidad y la veracidad. Si un científico publica un descubrimiento falso, tarde o temprano otros científicos intentarán replicar el experimento. Al no poder obtener los mismos resultados, el engaño quedará expuesto. 

\subsubsection*{Conclusión}
El castigo por el fraude científico es la excomunión profesional. El científico pierde toda credibilidad, el respeto de sus colegas y, en la inmensa mayoría de los casos, no vuelve a trabajar en el ámbito científico ni a publicar en revistas serias.

\newpage

\subsection*{Enunciado}
Contexto: Identificación de hipótesis científicas.
Pregunta: ¿Es la siguiente afirmación una hipótesis científica? (a) La clorofila hace verde al pasto.

\subsection*{Solución}

\subsubsection*{Definición de Hipótesis Científica}
Para que una afirmación sea considerada una hipótesis científica, no basta con que suene lógica o plausible. Según el principio de falsabilidad (central en la filosofía de la ciencia de Karl Popper y adoptado en la Física Conceptual de Hewitt), debe existir una prueba o experimento que sea capaz de demostrar que la hipótesis es errónea, en caso de que lo sea. Si no hay forma de someterla a una prueba de validez que pueda contradecirla, no es científica.

\subsubsection*{Análisis de la Afirmación}
La afirmación "La clorofila hace verde al pasto" vincula una sustancia química (la clorofila) con una propiedad observable (el color verde). Para determinar si es científica, debemos preguntarnos: ¿Existe un experimento que pueda demostrar que esto es falso? 
Sí existe. Podríamos extraer la clorofila de una muestra de pasto y observar si pierde su color verde, o podríamos cultivar pasto mutante que carezca de clorofila y verificar su color. Como la afirmación puede someterse a prueba mediante experimentación empírica, cumple con el criterio fundamental.

\subsubsection*{Conclusión}
Sí, es una hipótesis científica porque es comprobable y existe un experimento concebible que podría demostrar su falsedad si estuviera equivocada.

\newpage

\subsection*{Enunciado}
Contexto: Identificación de hipótesis científicas.
Pregunta: ¿Es la siguiente afirmación una hipótesis científica? (b) La Tierra gira en torno a su eje porque las cosas vivientes necesitan alternancia de luz y oscuridad.

\subsection*{Solución}

\subsubsection*{Causalidad Física vs. Teleología}
La física y las ciencias naturales se enfocan en los mecanismos y las causas físicas de los fenómenos ("cómo" y "por qué" en términos de fuerzas, energía y leyes naturales). La afirmación del problema es de naturaleza teleológica; es decir, atribuye un propósito o intención a un fenómeno natural (la rotación de la Tierra ocurre "para" satisfacer una necesidad de los seres vivos).

\subsubsection*{Prueba de Falsabilidad}
Aplicando el criterio de falsabilidad, debemos preguntarnos: ¿Qué experimento físico podríamos diseñar para demostrar que la Tierra NO gira por esta razón? No existe ninguno. Las necesidades de los seres vivos y la rotación planetaria son fenómenos que coexisten, pero atribuirle una intención cósmica al movimiento del planeta no es medible, cuantificable, ni sujeto a experimentación. Es una afirmación filosófica o especulativa, no científica.

\subsubsection*{Conclusión}
No es una hipótesis científica, ya que no puede ser sometida a pruebas experimentales para demostrar su falsedad; plantea un propósito en lugar de un mecanismo físico comprobable.

\newpage

\subsection*{Enunciado}
Contexto: Identificación de hipótesis científicas.
Pregunta: ¿Es la siguiente afirmación una hipótesis científica? (c) Las mareas son provocadas por la Luna.

\subsection*{Solución}

\subsubsection*{Relación de Variables Físicas}
Esta afirmación establece una relación de causa y efecto entre un cuerpo celeste (la Luna) y un fenómeno físico observable en la Tierra (las mareas oceánicas). Para ser considerada una hipótesis científica, esta afirmación debe predecir eventos que puedan ser medidos y verificados, y debe poder probarse si es falsa.

\subsubsection*{Análisis de Comprobación}
Podemos diseñar observaciones sistemáticas para poner a prueba esta idea. Si las mareas son provocadas por la Luna, entonces el ciclo de las mareas altas y bajas debería correlacionarse matemáticamente con la posición y las fases de la Luna respecto a la Tierra. Si, tras realizar mediciones exhaustivas, descubriéramos que los niveles del mar fluctúan de forma totalmente independiente de la trayectoria lunar, habríamos demostrado que la hipótesis es falsa. Como existe un método empírico para validarla o refutarla, cumple con los requisitos del método científico.

\subsubsection*{Conclusión}
Sí, es una hipótesis científica porque establece una causa física verificable y puede ser refutada mediante observaciones precisas de la órbita lunar y los niveles del océano.

\newpage

\subsection*{Enunciado}
Para responder a la pregunta de "cuando una planta crece, ¿de dónde proviene el material?", Aristóteles planteó por lógica la hipótesis de que todo el material provenía del suelo. ¿Consideras que su hipótesis es correcta, incorrecta o parcialmente correcta? ¿Qué experimentos propones para apoyar tu elección?

\subsection*{Solución}

\subsubsection*{Análisis de la Hipótesis Aristotélica}
El razonamiento de Aristóteles partía de la observación directa y la lógica pura: si una semilla se coloca en la tierra y de ella surge un gran árbol, el material sólido del árbol debía ser tierra transformada. Sin embargo, en la ciencia, la lógica sin experimentación puede llevar a conclusiones engañosas. Desde la perspectiva biológica y física moderna, sabemos que la inmensa mayoría de la masa de una planta (su celulosa y estructuras de carbono) proviene del dióxido de carbono del aire capturado durante la fotosíntesis, y del agua. El suelo solo aporta una cantidad ínfima de minerales y nutrientes.

\subsubsection*{Propuesta Experimental (Método Cuantitativo)}
Para refutar a Aristóteles científicamente, se requiere un experimento que aplique el principio de conservación de la masa. El experimento ideal es el históricamente realizado por el científico Jan van Helmont en el siglo XVII.
El procedimiento propuesto es el siguiente:
1. Secar completamente una gran cantidad de tierra y medir su masa con exactitud.
2. Colocar la tierra en una maceta y plantar un pequeño árbol joven cuya masa también haya sido medida exactamente.
3. Durante un periodo prolongado (por ejemplo, cinco años), regar la planta únicamente con agua pura (sin agregar tierra ni fertilizantes sólidos).
4. Al finalizar el periodo, retirar el árbol con cuidado, limpiar sus raíces y medir su nueva masa.
5. Volver a secar toda la tierra de la maceta y medir su masa final.

\subsubsection*{Justificación del Experimento}
Si la hipótesis de Aristóteles fuera correcta, la gran cantidad de masa ganada por el árbol debería ser exactamente igual a la cantidad de masa perdida por la tierra. Sin embargo, los resultados reales del experimento demostrarán que la tierra pierde apenas unos pocos gramos (los minerales), mientras que el árbol gana decenas de kilogramos. Esto demuestra de forma irrefutable que el material sólido proviene de otra fuente (el aire y el agua), falsando la lógica aristotélica.

\subsubsection*{Conclusión}
La hipótesis de Aristóteles es incorrecta. Para apoyar esta elección, propongo un experimento de balance de masas: medir la masa de la tierra y de la planta antes y después de varios años de crecimiento. Al evidenciar que la planta gana una enorme cantidad de masa mientras que el suelo casi no pierde nada, se demuestra que el material no proviene de la tierra.

\newpage

\subsection*{Enunciado}
¿Qué es lo que probablemente entiende mal una persona que dice, "pero esa sólo es una teoría científica"?

\subsection*{Solución}

\subsubsection*{El Significado Coloquial frente al Científico}
En el lenguaje cotidiano, la palabra "teoría" se utiliza a menudo como sinónimo de una suposición, una corazonada o una hipótesis que aún no ha sido comprobada. Cuando una persona no familiarizada con el método científico utiliza esta frase, asume que una "teoría científica" es simplemente una idea vacilante que carece de evidencia sólida.

\subsubsection*{La Definición en la Física Conceptual}
Según la epistemología expuesta por Paul G. Hewitt, en el ámbito de la ciencia, una "teoría" es exactamente lo contrario a una mera suposición. Una teoría científica es una síntesis profunda y estructurada de un gran cuerpo de información. Engloba hechos comprobados, leyes naturales y múltiples hipótesis que han sido rigurosamente probadas y verificadas por la comunidad científica a lo largo del tiempo. Ejemplos de esto son la Teoría de la Relatividad, la Teoría de la Evolución o la Teoría Atómica; marcos explicativos que sustentan nuestra comprensión moderna del universo.

\subsubsection*{Conclusión}
La persona entiende mal el rigor de la palabra. Confunde el uso coloquial (una conjetura sin pruebas) con la definición científica (un marco conceptual robusto, validado empíricamente y que sintetiza hechos y leyes naturales).

\newpage

\subsection*{Enunciado}
Se observa que la sombra proyectada por un pilar vertical en Alejandría al mediodía durante el solsticio de verano es $1/8$ la altura del pilar. La distancia entre Alejandría y Cirene es $1/8$ el radio de la Tierra. ¿Existe una conexión geométrica entre estas dos razones 1 a 8?

\subsection*{Solución}

\subsubsection*{Fundamentos del Experimento de Eratóstenes}
Este problema evoca la histórica medición de la Tierra realizada por Eratóstenes. En la ciudad de Cirene (Siena), durante el solsticio de verano al mediodía, los rayos del sol inciden de forma perpendicular a la superficie, sin proyectar sombra. En Alejandría, ubicada más al norte, los rayos solares paralelos inciden con un ángulo $\theta$ respecto a la vertical, proyectando una sombra. La geometría del sistema dicta que el ángulo de incidencia de los rayos solares en Alejandría es alterno interno e igual al ángulo $\theta$ que subtiende el arco entre las dos ciudades desde el centro de la Tierra.

\subsubsection*{Análisis Trigonométrico de la Sombra}
En el pilar de Alejandría se forma un triángulo rectángulo. El cateto opuesto al ángulo $\theta$ es la longitud de la sombra, denotada como $s$. El cateto adyacente es la altura del pilar, denotada como $h$. El problema establece la relación de la sombra con la altura, la cual corresponde a la función trigonométrica tangente:
$$\tan(\theta)=\frac{s}{h}$$
Sustituyendo el valor dado por el enunciado:
$$\tan(\theta)=\frac{1}{8}$$

\subsubsection*{Aproximación de Ángulos Pequeños}
En matemáticas, cuando un ángulo $\theta$ es suficientemente pequeño y se expresa en radianes, el valor de la tangente geométrica es prácticamente igual al valor del ángulo mismo. Es decir:
$$\tan(\theta)\approx\theta$$
Dado que $1/8$ es un valor pequeño (0.125), podemos concluir con un alto grado de precisión que el ángulo desde el centro de la Tierra equivale a $1/8$ de radián.
$$\theta\approx\frac{1}{8}\text{ radianes}$$

\subsubsection*{Relación de la Longitud de Arco}
La distancia lineal a lo largo de la superficie de una esfera entre dos puntos se conoce como longitud de arco, denotada como $D$. La fórmula que relaciona la longitud de arco, el radio $R$ y el ángulo central en radianes $\theta$ es:
$$D=R\cdot\theta$$
Sustituyendo el valor del ángulo que acabamos de deducir matemáticamente:
$$D=R\cdot\left(\frac{1}{8}\right)$$
Esto indica que la distancia lineal entre Alejandría y Cirene es, efectivamente, una octava parte del radio de la Tierra.

\subsubsection*{Representación Computacional de la Geometría}
Para consolidar la comprensión pedagógica de cómo el triángulo microscópico del pilar es homotético al sector circular macroscópico de la Tierra, se presenta el siguiente código que modela la geometría.

\begin{lstlisting}[language=Python]
import matplotlib.pyplot as plt
import numpy as np

fig, ax = plt.subplots(figsize=(8, 8))

theta_arc = np.linspace(np.pi/4, np.pi/2, 100)
x_earth = np.cos(theta_arc)
y_earth = np.sin(theta_arc)
ax.plot(x_earth, y_earth, color="blue", linewidth=2, label="Superficie Terrestre")

ax.plot([0, 0], [0, 1.5], color="gray", linestyle="--", label="Vertical en Cirene")
ax.plot([0, np.cos(1.2)], [0, np.sin(1.2)], color="black", linewidth=2, label="Pilar en Alejandria")

x_ray = np.linspace(0.8, 1.2, 10)
y_ray = np.linspace(1.5, 1.5, 10)
for i in range(3):
    ax.plot(x_ray, y_ray - i*0.2, color="orange", label="Rayos Solares" if i==0 else "")

ax.set_aspect("equal")
plt.legend()
plt.axis("off")
plt.show()
\end{lstlisting}

\subsubsection*{Conclusión}
Sí, existe una conexión geométrica directa y causal. El ángulo de incidencia de los rayos solares (cuya tangente es $1/8$) corresponde en radianes aproximadamente a $1/8$. Puesto que el arco terrestre se define como $D=R\cdot\theta$, al multiplicar el radio de la Tierra por este ángulo en radianes, se obtiene lógicamente que la distancia entre las ciudades es $1/8$ del radio terrestre.

\newpage

\subsection*{Enunciado}
Si la Tierra fuera más pequeña de lo que es, pero la distancia de Alejandría a Cirene fuera la misma, ¿la sombra del pilar vertical en Alejandría sería más larga o más corta al mediodía durante el solsticio de verano?

\subsection*{Solución}

\subsubsection*{Análisis de la Curvatura y el Ángulo Central}
Partimos de la relación geométrica establecida en el problema anterior, donde la distancia $D$ (longitud de arco entre Alejandría y Cirene) está dada por $D=R\cdot\theta$.
Si despejamos el ángulo central $\theta$, obtenemos:
$$\theta=\frac{D}{R}$$
El problema plantea un escenario hipotético donde la distancia $D$ se mantiene constante, pero el planeta es más pequeño, lo que significa que el radio $R$ disminuye. Matemáticamente, al dividir una constante ($D$) por un denominador más pequeño ($R$), el resultado del cociente es mayor. Por lo tanto, el ángulo central $\theta$ aumentaría.
Físicamente, esto significa que un planeta más pequeño tiene una curvatura más pronunciada, provocando que la vertical en Alejandría se incline de forma más abrupta respecto a la vertical en Cirene.

\subsubsection*{Efecto del Ángulo Aumentado en la Sombra}
Como sabemos, el ángulo de incidencia de los rayos solares respecto al pilar en Alejandría es geométricamente idéntico a este ángulo central $\theta$. 
La longitud de la sombra $s$ se determina despejando la ecuación trigonométrica:
$$s=h\cdot\tan(\theta)$$
Donde $h$ (la altura del pilar) permanece constante. En el primer cuadrante de las funciones trigonométricas, a medida que un ángulo $\theta$ aumenta, su tangente también aumenta. Dado que la nueva Tierra más pequeña genera un ángulo $\theta$ mayor, el valor de $\tan(\theta)$ será mayor. Al multiplicar esto por la altura del pilar, obtendremos obligatoriamente un valor de $s$ más grande.

\subsubsection*{Representación Computacional de la Curvatura}
El siguiente código permite visualizar cómo un mismo arco proyecta un ángulo mayor sobre una circunferencia de menor radio, impactando la inclinación del pilar.

\begin{lstlisting}[language=Python]
import matplotlib.pyplot as plt
import numpy as np

fig, ax = plt.subplots(figsize=(6, 6))

theta_large = np.linspace(0, np.pi/2, 100)
ax.plot(2 * np.cos(theta_large), 2 * np.sin(theta_large), color="blue", label="Tierra Original")

theta_small = np.linspace(0, np.pi, 100)
ax.plot(1 * np.cos(theta_small), 1 * np.sin(theta_small), color="red", label="Tierra Mas Pequena")

ax.set_aspect("equal")
plt.legend()
plt.axis("off")
plt.show()
\end{lstlisting}

\subsubsection*{Conclusión}
La sombra del pilar sería más larga. Al disminuir el radio de la Tierra manteniendo la misma distancia entre las ciudades, la curvatura entre ellas es más aguda. Esto incrementa el ángulo con el que los rayos solares inciden sobre la vertical del pilar, proyectando así una sombra de mayor longitud.


\newpage

\subsection*{Enunciado}
Lucy Piesligeros está parada con un pie en una báscula de baño y el otro pie en una segunda báscula de baño. La lectura de cada báscula es de $350\text{ N}$. ¿Cuál es el peso de Lucy?

\subsection*{Solución}

\subsubsection*{El Principio de Equilibrio Estático}
De acuerdo con las leyes de Newton abordadas en la Física Conceptual, cuando un objeto físico se encuentra en reposo absoluto respecto a su marco de referencia, está en un estado de equilibrio estático. En este estado, la regla del equilibrio dicta que la suma vectorial de todas las fuerzas (la fuerza neta) que actúan sobre el objeto debe ser igual a cero. Matemáticamente, esto se expresa mediante la ecuación $\Sigma F=0$.

\subsubsection*{Análisis de las Fuerzas Intervinientes}
Sobre el cuerpo de Lucy actúan simultáneamente fuerzas colineales en el eje vertical. Hacia abajo actúa la fuerza de atracción gravitacional de la Tierra, que constituye su peso. Hacia arriba actúan las fuerzas de soporte proporcionadas por las superficies de las básculas. Es importante aclarar que una báscula no mide el peso directamente, sino la fuerza de soporte (fuerza normal) que la propia báscula debe ejercer hacia arriba contra el pie para contrarrestar la gravedad. Como Lucy tiene un pie en cada báscula, existen dos fuerzas normales independientes dirigidas hacia arriba empujándola.

\subsubsection*{Diagrama de Cuerpo Libre Computacional}
Para comprender pedagógicamente el equilibrio mecánico, se presenta el siguiente script que modela el diagrama de cuerpo libre. Se observa cómo la suma geométrica de los vectores ascendentes de ambas básculas debe equilibrar y ser idéntica en magnitud al vector descendente del peso.

\begin{lstlisting}[language=Python]
import matplotlib.pyplot as plt

fig, ax = plt.subplots(figsize=(6, 7))

ax.plot(0, 0, 'ko', markersize=12)

ax.annotate("", xy=(0, -700), xytext=(0, 0), arrowprops=dict(facecolor='red', width=4, headwidth=12))
ax.text(0.2, -350, "W = ?", color="red", fontsize=14, fontweight='bold')

ax.annotate("", xy=(-0.5, 350), xytext=(0, 0), arrowprops=dict(facecolor='blue', width=4, headwidth=12))
ax.text(-1.8, 175, "N1 = 350 N", color="blue", fontsize=12)

ax.annotate("", xy=(0.5, 350), xytext=(0, 0), arrowprops=dict(facecolor='green', width=4, headwidth=12))
ax.text(0.7, 175, "N2 = 350 N", color="green", fontsize=12)

ax.set_xlim(-2.5, 2.5)
ax.set_ylim(-800, 500)
plt.title("Diagrama de Cuerpo Libre: Equilibrio de Fuerzas")
plt.axis("off")
plt.show()
\end{lstlisting}

\subsubsection*{Cálculo Matemático}
Aplicando la primera condición de equilibrio exclusivamente en el eje vertical, la suma de las fuerzas que apuntan hacia arriba debe ser igual a la magnitud de la fuerza que apunta hacia abajo.
$$\Sigma F_{y}=0$$
$$N_{1}+N_{2}-W=0$$
Al transponer el término del peso $W$ al otro lado de la ecuación obtenemos:
$$W=N_{1}+N_{2}$$
Sustituyendo los valores numéricos de las lecturas de las básculas proporcionados en el enunciado:
$$W=350\text{ N}+350\text{ N}$$
$$W=700\text{ N}$$

\subsubsection*{Conclusión}
El peso total de Lucy es de $700\text{ N}$.


\newpage

\subsection*{Enunciado}
Henry Pesopesado pesa $1,200\text{ N}$ y se para sobre un par de básculas de baño de modo que una báscula lee el doble que la otra. ¿Cuáles son las lecturas de las básculas?

\subsection*{Solución}

\subsubsection*{El Principio de Equilibrio Mecánico}
Al igual que en el ejercicio anterior, nos basamos en la primera ley de Newton y el concepto de equilibrio estático presentado en la Física Conceptual de Hewitt. Dado que Henry se encuentra en reposo sobre las básculas, la fuerza neta que actúa sobre él es cero. Esto significa que las fuerzas descendentes (su peso debido a la gravedad) deben ser perfectamente contrarrestadas por las fuerzas ascendentes (las fuerzas normales o de soporte ejercidas por ambas básculas).

\subsubsection*{Modelado Algebraico de las Fuerzas}
Definiremos las variables de nuestro sistema. Sea $W$ el peso total de Henry, con un valor conocido de $1,200\text{ N}$. Sean $N_{1}$ y $N_{2}$ las fuerzas de soporte individuales proporcionadas por cada báscula. 
La ecuación de equilibrio en el eje vertical se expresa de la siguiente manera:
$$\Sigma F_{y}=0$$
$$N_{1}+N_{2}-W=0$$
$$N_{1}+N_{2}=1,200\text{ N}$$

\subsubsection*{Condición de Proporcionalidad}
El problema introduce una restricción asimétrica: Henry no está distribuyendo su peso equitativamente. Se nos indica que una báscula lee el doble que la otra. Matemáticamente, esto establece un sistema de ecuaciones donde una variable es dependiente de la otra mediante un factor de escala. Podemos expresar esta relación como:
$$N_{1}=2\cdot N_{2}$$

\subsubsection*{Diagrama de Cuerpo Libre Computacional}
Para visualizar cómo se distribuye la fuerza de forma desigual pero manteniendo el equilibrio estático total, se presenta el siguiente código. Observamos cómo la suma de los vectores ascendentes asimétricos equilibra al gran vector descendente.

\begin{lstlisting}[language=Python]
import matplotlib.pyplot as plt

fig, ax = plt.subplots(figsize=(6, 8))

ax.plot(0, 0, 'ko', markersize=12)

ax.annotate("", xy=(0, -1200), xytext=(0, 0), arrowprops=dict(facecolor='red', width=4, headwidth=12))
ax.text(0.2, -600, "W = 1200 N", color="red", fontsize=12, fontweight='bold')

ax.annotate("", xy=(-0.5, 800), xytext=(0, 0), arrowprops=dict(facecolor='blue', width=4, headwidth=12))
ax.text(-1.8, 400, "N1 = 2x", color="blue", fontsize=12)

ax.annotate("", xy=(0.5, 400), xytext=(0, 0), arrowprops=dict(facecolor='green', width=4, headwidth=12))
ax.text(0.7, 200, "N2 = x", color="green", fontsize=12)

ax.set_xlim(-2.5, 2.5)
ax.set_ylim(-1300, 900)
plt.title("Equilibrio Estatico con Carga Asimetrica")
plt.axis("off")
plt.show()
\end{lstlisting}

\subsubsection*{Resolución del Sistema de Ecuaciones}
Sustituimos la condición de proporcionalidad en la ecuación de equilibrio principal. Donde dice $N_{1}$, colocaremos su equivalente $2\cdot N_{2}$:
$$(2\cdot N_{2})+N_{2}=1,200\text{ N}$$
Sumamos los términos semejantes:
$$3\cdot N_{2}=1,200\text{ N}$$
Despejamos $N_{2}$ dividiendo ambos lados entre 3:
$$N_{2}=\frac{1,200\text{ N}}{3}$$
$$N_{2}=400\text{ N}$$
Ahora que conocemos el valor de $N_{2}$, utilizamos la condición de proporcionalidad para encontrar $N_{1}$:
$$N_{1}=2\cdot(400\text{ N})$$
$$N_{1}=800\text{ N}$$
Para comprobar, sumamos ambas lecturas: $800\text{ N} + 400\text{ N} = 1,200\text{ N}$, lo cual satisface la condición de equilibrio estático.

\subsubsection*{Conclusión}
Las lecturas de las básculas son de $800\text{ N}$ y $400\text{ N}$ respectivamente.

\newpage

\subsection*{Enunciado}
El dibujo muestra el andamio de un pintor en equilibrio mecánico. La persona que está en medio del andamio pesa $500\text{ N}$ y las tensiones en cada soga son de $400\text{ N}$. ¿Cuál es el peso del andamio?

\begin{center}
\begin{tikzpicture}[scale=1]
  \draw[thick, fill=brown!40] (-3, 0) rectangle (3, 0.2);
  \draw[thick] (-2.8, 0.2) -- (-2.8, 2.5);
  \draw[thick] (2.8, 0.2) -- (2.8, 2.5);
  
  \draw[->, >=stealth, very thick, magenta] (-2.8, 0) -- (-2.8, 1.5) node[right, black] {400 N};
  \draw[->, >=stealth, very thick, magenta] (2.8, 0) -- (2.8, 1.5) node[left, black] {400 N};
  
  \draw[thick, fill=blue!20] (-0.2, 0.2) rectangle (0.2, 1);
  \draw[thick] (0, 1.2) circle (0.2);
  
  \draw[->, >=stealth, very thick, magenta] (0, 0.2) -- (0, -1.5) node[left, black] {500 N};
\end{tikzpicture}
\end{center}

\subsection*{Solución}

\subsubsection*{Definición del Sistema y Equilibrio Mecánico}
Para resolver este problema, debemos considerar todo el conjunto (la persona y el andamio) como un único sistema físico. Según los principios expuestos en la Física Conceptual, cuando un sistema está en equilibrio mecánico y no presenta aceleración, la regla del equilibrio establece que la fuerza neta es igual a cero. En este caso unidimensional (fuerzas verticales), la suma de las fuerzas que tiran hacia arriba debe ser exactamente igual a la suma de las fuerzas que tiran hacia abajo.

\subsubsection*{Identificación de las Fuerzas}
En este sistema, interactúan cuatro fuerzas verticales:
Hacia arriba actúan las fuerzas de tensión de las dos sogas que sostienen el andamio. Llamémoslas $T_{1}$ y $T_{2}$, cada una con un valor de $400\text{ N}$.
Hacia abajo actúan las fuerzas gravitacionales (los pesos). Por un lado, tenemos el peso de la persona, $W_{p}$, que es de $500\text{ N}$. Por otro lado, tenemos el peso del propio andamio, $W_{a}$, que es el valor desconocido que debemos calcular.

\subsubsection*{Diagrama de Cuerpo Libre Computacional}
El siguiente código modela las fuerzas que actúan sobre el centro de masa del sistema. Se puede apreciar visualmente cómo los dos vectores de tensión ascendentes deben compensar tanto al vector del peso del pintor como al vector del peso del andamio para mantener el equilibrio.

\begin{lstlisting}[language=Python]
import matplotlib.pyplot as plt

fig, ax = plt.subplots(figsize=(6, 7))

ax.plot(0, 0, 'ko', markersize=10)

ax.annotate("", xy=(-0.5, 400), xytext=(0, 0), arrowprops=dict(facecolor='magenta', width=3, headwidth=10))
ax.text(-1.8, 200, "T1 = 400 N", color="magenta", fontsize=11, fontweight='bold')

ax.annotate("", xy=(0.5, 400), xytext=(0, 0), arrowprops=dict(facecolor='magenta', width=3, headwidth=10))
ax.text(0.7, 200, "T2 = 400 N", color="magenta", fontsize=11, fontweight='bold')

ax.annotate("", xy=(0, -500), xytext=(0, 0), arrowprops=dict(facecolor='blue', width=3, headwidth=10))
ax.text(0.2, -250, "W_persona = 500 N", color="blue", fontsize=11)

ax.annotate("", xy=(0, -800), xytext=(0, -500), arrowprops=dict(facecolor='orange', width=3, headwidth=10))
ax.text(0.2, -650, "W_andamio = ?", color="orange", fontsize=11)

ax.set_xlim(-2.5, 2.5)
ax.set_ylim(-900, 500)
plt.title("DCL: Equilibrio de Tensiones y Pesos")
plt.axis("off")
plt.show()
\end{lstlisting}

\subsubsection*{Planteamiento Matemático y Resolución}
Aplicamos la primera condición de equilibrio estableciendo que la sumatoria de fuerzas en el eje de las ordenadas es igual a cero:
$$\Sigma F_{y}=0$$
Consideramos positivas las fuerzas dirigidas hacia arriba y negativas las dirigidas hacia abajo:
$$T_{1}+T_{2}-W_{p}-W_{a}=0$$
Sustituimos los valores numéricos proporcionados por el problema:
$$400\text{ N}+400\text{ N}-500\text{ N}-W_{a}=0$$
Sumamos las fuerzas de tensión:
$$800\text{ N}-500\text{ N}-W_{a}=0$$
Restamos el peso de la persona de la fuerza ascendente total:
$$300\text{ N}-W_{a}=0$$
Finalmente, despejamos la variable desconocida transponiéndola al otro lado de la igualdad:
$$W_{a}=300\text{ N}$$

\subsubsection*{Conclusión}
El peso del andamio es de $300\text{ N}$.


\newpage
\subsection*{Enunciado}
Un andamio diferente, que pesa 400 N, soporta dos pintores, uno de 500 N y otro de 400 N. La lectura de la báscula izquierda es de 800 N. ¿Cuál es la lectura de la báscula de la derecha?

\begin{center}
\begin{tikzpicture}[scale=1]
  \draw[thick, fill=brown!40] (-4, 0) rectangle (4, 0.2);
  
  \draw[thick] (-3.8, 0.2) -- (-3.8, 3.5);
  \draw[thick] (3.8, 0.2) -- (3.8, 3.5);
  
  \draw[thick, fill=purple!50] (-3.5, 0.2) rectangle (-2.9, 1);
  \draw[thick] (-3.2, 1.2) circle (0.2);
  \draw[thick] (-3.7, 0.6) -- (-3.5, 0.4);
  \draw[thick] (-2.7, 0.6) -- (-2.9, 0.4);
  
  \draw[thick, fill=green!50] (2.5, 0.2) rectangle (3.1, 0.9);
  \draw[thick] (2.8, 1.1) circle (0.2);
  
  \draw[->, >=stealth, very thick, magenta] (-3.8, 0.2) -- (-3.8, 3) node[right, black] {800 N};
  \node[right, font=\large] at (3.9, 1.8) {?};
  
  \draw[->, >=stealth, very thick, magenta] (-3.2, 0) -- (-3.2, -1.8) node[right, black] {500 N};
  
  \draw[->, >=stealth, very thick, magenta] (0, 0) -- (0, -1.3) node[right, black] {400 N};
  
  \draw[->, >=stealth, very thick, magenta] (2.8, 0) -- (2.8, -1.3) node[right, black] {400 N};
\end{tikzpicture}
\end{center}

\subsection*{Solución}

\subsubsection*{El Principio de Equilibrio Traslacional}
Para abordar este problema desde la perspectiva de la Física Conceptual, definimos a todo el conjunto (el andamio y los dos pintores) como un único sistema físico en estado de reposo. Cuando un sistema no experimenta cambios en su estado de movimiento, se encuentra en equilibrio mecánico. La condición fundamental para el equilibrio traslacional es que la suma vectorial de todas las fuerzas externas que actúan sobre el sistema sea igual a cero.

\subsubsection*{Identificación y Clasificación de las Fuerzas}
En este escenario unidimensional, todas las fuerzas actúan en el eje vertical. Podemos clasificarlas en dos grupos:
Fuerzas ascendentes: Son las tensiones generadas por las sogas que sostienen el andamio. Tenemos la tensión izquierda conocida de 800 N y la tensión derecha desconocida, que llamaremos $T_{d}$.
Fuerzas descendentes: Corresponden a los pesos ocasionados por la atracción gravitatoria. Tenemos el peso del primer pintor de 500 N, el peso del propio andamio de 400 N y el peso del segundo pintor de 400 N.

\subsubsection*{Diagrama de Cuerpo Libre Computacional}
Para consolidar visualmente la regla del equilibrio, el siguiente script genera un diagrama de cuerpo libre donde se agrupan las fuerzas. Se demuestra que la magnitud total de los vectores que apuntan hacia arriba debe ser geométricamente idéntica a la suma de los vectores que apuntan hacia abajo.

\begin{lstlisting}[language=Python]
import matplotlib.pyplot as plt

fig, ax = plt.subplots(figsize=(7, 8))

ax.plot(0, 0, 'ko', markersize=12)

ax.annotate("", xy=(-0.5, 800), xytext=(0, 0), arrowprops=dict(facecolor='magenta', width=3, headwidth=10))
ax.text(-1.8, 400, "T_izq = 800 N", color="magenta", fontsize=11, fontweight='bold')

ax.annotate("", xy=(0.5, 500), xytext=(0, 0), arrowprops=dict(facecolor='orange', width=3, headwidth=10))
ax.text(0.7, 250, "T_der = ?", color="orange", fontsize=11, fontweight='bold')

ax.annotate("", xy=(-0.5, -500), xytext=(0, 0), arrowprops=dict(facecolor='blue', width=3, headwidth=10))
ax.text(-2.5, -250, "W_pintor1 = 500 N", color="blue", fontsize=10)

ax.annotate("", xy=(0, -400), xytext=(0, 0), arrowprops=dict(facecolor='green', width=3, headwidth=10))
ax.text(0.2, -200, "W_andamio = 400 N", color="green", fontsize=10)

ax.annotate("", xy=(0.5, -400), xytext=(0, 0), arrowprops=dict(facecolor='cyan', width=3, headwidth=10))
ax.text(0.7, -350, "W_pintor2 = 400 N", color="cyan", fontsize=10)

ax.set_xlim(-3, 3)
ax.set_ylim(-600, 1000)
plt.title("DCL: Equilibrio del Sistema Completo")
plt.axis("off")
plt.show()
\end{lstlisting}

\subsubsection*{Planteamiento Matemático y Resolución}
Aplicamos la primera ley de Newton para el eje vertical, estableciendo que la sumatoria de fuerzas es nula:
$$\Sigma F_{y}=0$$
Asignamos signos positivos a las fuerzas dirigidas hacia arriba y negativos a las fuerzas dirigidas hacia abajo:
$$T_{i}+T_{d}-W_{p1}-W_{a}-W_{p2}=0$$
Sustituimos los valores numéricos proporcionados en el enunciado:
$$800+T_{d}-500-400-400=0$$
Agrupamos y sumamos todos los valores correspondientes a los pesos (fuerzas descendentes):
$$800+T_{d}-1300=0$$
Realizamos la operación aritmética entre las constantes en el lado izquierdo de la ecuación:
$$T_{d}-500=0$$
Finalmente, despejamos la variable de la tensión derecha aislándola en un lado de la igualdad:
$$T_{d}=500$$

\subsubsection*{Conclusión}
La lectura de la báscula de la derecha es de 500 N.

\newpage

\subsection*{Enunciado}
Los pesos de Burl, Paul y el andamio producen tensiones en las sogas de sostén. Clasifica las tensiones en la soga \textit{izquierda}, de mayor a menor, en las tres situaciones, A, B y C.

\vspace{0.5cm}
\begin{center}
\begin{tikzpicture}[scale=0.8]
  \draw[thick, fill=brown!40] (-3, 0) rectangle (2, 0.2);
  \draw[thick] (-2.5, 0.2) -- (-2.5, 3);
  \draw[thick] (1.5, 0.2) -- (1.5, 3);
  \draw[thick, fill=red!60] (-2.7, 0.2) rectangle (-2.1, 1);
  \draw[thick] (-2.4, 1.2) circle (0.2);
  \draw[thick, fill=blue!50] (1.2, 0.2) rectangle (1.7, 0.9);
  \draw[thick] (1.45, 1.05) circle (0.15);
  \node[font=\large\bfseries] at (-0.5, -0.5) {A};

  \begin{scope}[xshift=6cm]
    \draw[thick, fill=brown!40] (-3, 0) rectangle (2, 0.2);
    \draw[thick] (-2.5, 0.2) -- (-2.5, 3);
    \draw[thick] (1.5, 0.2) -- (1.5, 3);
    \draw[thick, fill=red!60] (-2.7, 0.2) rectangle (-2.1, 1);
    \draw[thick] (-2.4, 1.2) circle (0.2);
    \draw[thick, fill=blue!50] (-0.5, 0.2) rectangle (0, 0.9);
    \draw[thick] (-0.25, 1.05) circle (0.15);
    \node[font=\large\bfseries] at (-0.5, -0.5) {B};
  \end{scope}

  \begin{scope}[xshift=12cm]
    \draw[thick, fill=brown!40] (-3, 0) rectangle (2, 0.2);
    \draw[thick] (-2.5, 0.2) -- (-2.5, 3);
    \draw[thick] (1.5, 0.2) -- (1.5, 3);
    \draw[thick, fill=red!60] (-2.7, 0.2) rectangle (-2.1, 1);
    \draw[thick] (-2.4, 1.2) circle (0.2);
    \draw[thick, fill=blue!50] (-1.5, 0.2) rectangle (-1.0, 0.9);
    \draw[thick] (-1.25, 1.05) circle (0.15);
    \node[font=\large\bfseries] at (-0.5, -0.5) {C};
  \end{scope}
\end{tikzpicture}
\end{center}

\subsection*{Solución}

\subsubsection*{Principios de Equilibrio y Momentos de Torsión}
De acuerdo con los conceptos de la Física Conceptual de Paul G. Hewitt, un objeto rígido (como el andamio) en equilibrio estático debe satisfacer dos condiciones. La primera es que la suma de fuerzas debe ser cero ($\Sigma F = 0$), lo que implica que la suma de las tensiones hacia arriba en ambas sogas es siempre igual al peso total hacia abajo (Burl + Paul + andamio). Como nadie se baja del andamio, la fuerza total hacia arriba que deben proveer las sogas es constante en las tres situaciones.
La segunda condición es que la suma de los momentos de torsión (torques) debe ser cero ($\Sigma \tau = 0$). El momento de torsión determina cómo se distribuye esa fuerza total entre los dos soportes. Cuanto más cerca esté el centro de gravedad del sistema de uno de los soportes, mayor será la proporción de peso que ese soporte deberá cargar.

\subsubsection*{Análisis de la Distribución del Peso}
Evaluemos qué ocurre en cada situación respecto a la soga izquierda:
El peso del andamio está en el centro, por lo que su carga se distribuye equitativamente entre la soga izquierda y la derecha en todo momento.
Burl permanece estático en el extremo izquierdo en las tres situaciones, por lo que la carga que él aporta a la soga izquierda es máxima y constante en A, B y C.
La variable que cambia es la posición de Paul. Como Paul se desplaza, cambia el centro de gravedad del sistema completo.

Situación A: Paul está en el extremo derecho. La soga derecha soporta casi todo el peso de Paul. La soga izquierda soporta muy poco del peso de Paul.
Situación B: Paul camina hacia el centro. Ahora, su peso se distribuye de manera más equilibrada entre ambas sogas. La tensión en la soga izquierda aumenta porque ahora carga una mayor fracción del peso de Paul que en la situación A.
Situación C: Paul se acerca a Burl, en el lado izquierdo. En esta posición, la inmensa mayoría de su peso es soportada por la soga izquierda.

\subsubsection*{Representación Computacional de las Tensiones}
El siguiente código modela el comportamiento de las fuerzas de soporte al mover una masa a lo largo de una viga, demostrando matemáticamente cómo la tensión izquierda alcanza su máximo cuando las masas se concentran en ese lado.

\begin{lstlisting}[language=Python]
import matplotlib.pyplot as plt
import numpy as np

W_andamio = 400
W_burl = 600
W_paul = 400
L = 10 

posiciones_paul = {'A': 10, 'B': 5, 'C': 2}
tensiones_izq = []

fig, axs = plt.subplots(1, 3, figsize=(12, 4), sharey=True)
fig.suptitle("Evolucion de Tensiones al Desplazar el Peso", fontsize=14)

for i, (situacion, x_paul) in enumerate(posiciones_paul.items()):
    sum_tau_derecha = (W_burl * 10) + (W_andamio * 5) + (W_paul * (10 - x_paul))
    T_izq = sum_tau_derecha / 10
    tensiones_izq.append(T_izq)
    T_der = (W_burl + W_andamio + W_paul) - T_izq
    
    axs[i].plot([0, L], [0, 0], color='brown', linewidth=6)
    axs[i].annotate("", xy=(0, T_izq), xytext=(0, 0), arrowprops=dict(facecolor='magenta', width=4))
    axs[i].annotate("", xy=(L, T_der), xytext=(L, 0), arrowprops=dict(facecolor='orange', width=4))
    axs[i].plot(0, -W_burl/10, 'ro', markersize=15, label="Burl")
    axs[i].plot(x_paul, -W_paul/10, 'bo', markersize=10, label="Paul")
    
    axs[i].set_title(f"Situacion {situacion}\nT_izq = {T_izq:.0f} N")
    axs[i].set_xlim(-1, 11)
    axs[i].set_ylim(-150, 1500)
    axs[i].axis('off')

plt.tight_layout()
plt.show()
\end{lstlisting}

\subsubsection*{Conclusión}
Dado que la tensión de la soga izquierda aumenta a medida que el peso total se acerca a ella, la clasificación de mayor a menor tensión es: C, B, A.


\newpage

\subsection*{Enunciado}
La soga sostiene una linterna que pesa $50\text{ N}$. ¿La tensión en la soga es menor que, igual a o mayor que $50\text{ N}$? Usa la regla del paralelogramo para defender tu respuesta.

\begin{center}
\begin{tikzpicture}[scale=1]
  \draw[thick] (-2, 2) -- (2, 2);
  \foreach \x in {-1.8, -1.4, -1.0, -0.6, -0.2, 0.2, 0.6, 1.0, 1.4, 1.8} 
    \draw (\x, 2) -- (\x+0.2, 2.2);
  \draw[thick] (-0.6, 2) -- (0, 0.5);
  \draw[thick] (0.6, 2) -- (0, 0.5);
  
  \draw[thick, fill=cyan!40] (-0.5, 0.5) rectangle (0.5, 0.2);
  \draw[thick, fill=yellow!60] (-0.4, 0.2) rectangle (0.4, -0.4);
  \draw[thick] (-0.15, 0.2) -- (-0.15, -0.4);
  \draw[thick] (0.15, 0.2) -- (0.15, -0.4);
  \draw[thick, fill=cyan!40] (-0.5, -0.4) rectangle (0.5, -0.7);
  \draw[thick] (0, 0.6) circle (0.1);
  \draw[thick] (0, 0.5) -- (0, 0.2);
\end{tikzpicture}
\end{center}

\subsection*{Solución}

\subsubsection*{Equilibrio y la Fuerza Resultante}
Para que la linterna permanezca colgada en reposo, debe cumplirse la primera condición de equilibrio mecánico propuesta en la Física Conceptual: la fuerza neta debe ser nula ($\Sigma F = 0$). Dado que la gravedad ejerce una fuerza de $50\text{ N}$ hacia abajo (el peso de la linterna), las sogas deben ejercer, en conjunto, una fuerza resultante exactamente igual a $50\text{ N}$ dirigida hacia arriba.

\subsubsection*{Aplicación de la Regla del Paralelogramo}
La fuerza ascendente total de $50\text{ N}$ es la diagonal del paralelogramo formado por los dos vectores de tensión de la soga. Como el ángulo que forman las dos mitades de la soga es relativamente pequeño (las sogas están casi verticales), el paralelogramo resultante es alto y estrecho. 
Geométricamente, cuando el ángulo entre dos vectores es pequeño y su diagonal resultante debe alcanzar un valor fijo (en este caso $50\text{ N}$), la longitud de los lados del paralelogramo (que representan las magnitudes de las tensiones) será menor que la longitud de la diagonal.

\subsubsection*{Demostración Computacional del Paralelogramo}
El siguiente script construye a escala el vector resultante necesario para mantener el equilibrio y grafica cómo se descompone en las dos sogas con un ángulo cerrado, demostrando visualmente que los vectores laterales son más cortos que la diagonal central.

\begin{lstlisting}[language=Python]
import matplotlib.pyplot as plt

fig, ax = plt.subplots(figsize=(4, 6))

ax.annotate("", xy=(0, 50), xytext=(0, 0), arrowprops=dict(facecolor='blue', width=3, headwidth=10))
ax.text(2, 25, "Resultante = 50 N", color="blue", fontsize=10, fontweight='bold')

ax.annotate("", xy=(-12, 25), xytext=(0, 0), arrowprops=dict(facecolor='red', width=2, headwidth=8))
ax.annotate("", xy=(12, 25), xytext=(0, 0), arrowprops=dict(facecolor='red', width=2, headwidth=8))
ax.text(-20, 12, "Tension", color="red", fontsize=10)
ax.text(14, 12, "Tension", color="red", fontsize=10)

ax.plot([-12, 0], [25, 50], color='gray', linestyle='--')
ax.plot([12, 0], [25, 50], color='gray', linestyle='--')

ax.set_xlim(-25, 25)
ax.set_ylim(-5, 55)
plt.title("Regla del Paralelogramo: Angulo Pequeno")
plt.axis("off")
plt.show()
\end{lstlisting}

\subsubsection*{Conclusión}
La tensión en la soga es menor que $50\text{ N}$. Al aplicar la regla del paralelogramo con un ángulo agudo entre las cuerdas, la longitud de los lados (tensión) resulta ser geométricamente inferior a la longitud de la diagonal (peso).

\newpage

\subsection*{Enunciado}
Un mono cuelga inmóvil al final de una liana vertical. ¿Cuáles dos fuerzas actúan sobre el mono? Si una de ellas es mayor, ¿cuál es?

\subsection*{Solución}

\subsubsection*{Identificación de Fuerzas}
Cualquier objeto masivo que se encuentre dentro del campo gravitatorio terrestre experimenta una fuerza de atracción hacia el centro del planeta; esta es la primera fuerza, denominada peso (o fuerza de gravedad), que actúa con dirección vertical hacia abajo. Por otro lado, la liana se estira bajo la masa del mono debido a sus propiedades elásticas a nivel molecular, generando una fuerza de restitución que tira del mono hacia arriba; esta es la segunda fuerza, denominada tensión.

\subsubsection*{Condición de Equilibrio Estático}
El término clave en el enunciado es "inmóvil". En términos físicos, esto implica que la aceleración del mono es nula ($a = 0$). De acuerdo con la Segunda Ley de Newton ($F_{neta} = m \cdot a$), si la aceleración es nula, la fuerza neta que actúa sobre el cuerpo debe ser obligatoriamente cero. Dado que solo existen dos fuerzas operando en la misma línea de acción (el eje vertical), la única forma matemática de que su suma sea cero es que ambas fuerzas posean exactamente la misma magnitud pero sentidos opuestos.

\subsubsection*{Conclusión}
Las dos fuerzas que actúan sobre el mono son el peso (hacia abajo) y la tensión de la liana (hacia arriba). Ninguna de las dos es mayor; deben ser exactamente iguales en magnitud para que el mono permanezca inmóvil en equilibrio estático.

\newpage
\subsection*{Enunciado}
La soga del ejercicio 56 se coloca como se muestra y aun así sostiene la linterna de $50\text{ N}$. ¿La tensión en la soga es menor que, igual a o mayor que $50\text{ N}$? Usa la regla del paralelogramo para defender tu respuesta.

\begin{center}
\begin{tikzpicture}[scale=1]
  \draw[thick] (-4, 2) -- (4, 2);
  \foreach \x in {-3.8, -3.4, -3.0, -2.6, -2.2, -1.8, -1.4, -1.0, -0.6, -0.2, 0.2, 0.6, 1.0, 1.4, 1.8, 2.2, 2.6, 3.0, 3.4, 3.8} 
    \draw (\x, 2) -- (\x+0.2, 2.2);
  \draw[thick] (-3, 2) -- (0, 0.5);
  \draw[thick] (3, 2) -- (0, 0.5);
  
  \draw[thick, fill=cyan!40] (-0.5, 0.5) rectangle (0.5, 0.2);
  \draw[thick, fill=yellow!60] (-0.4, 0.2) rectangle (0.4, -0.4);
  \draw[thick] (-0.15, 0.2) -- (-0.15, -0.4);
  \draw[thick] (0.15, 0.2) -- (0.15, -0.4);
  \draw[thick, fill=cyan!40] (-0.5, -0.4) rectangle (0.5, -0.7);
  \draw[thick] (0, 0.6) circle (0.1);
  \draw[thick] (0, 0.5) -- (0, 0.2);
\end{tikzpicture}
\end{center}

\subsection*{Solución}

\subsubsection*{Equilibrio con Ángulo Obtuso}
El principio físico inicial es el mismo que en el primer ejercicio: la diagonal vertical resultante (fuerza neta hacia arriba) debe ser de exactamente $50\text{ N}$ para equilibrar el peso de la linterna. Sin embargo, en esta configuración, la geometría del sistema de soporte ha cambiado drásticamente. Las sogas forman un ángulo muy abierto (obtuso) entre sí, acercándose a la horizontal.

\subsubsection*{Geometría Extrema del Paralelogramo}
Cuando dibujamos el vector vertical resultante de $50\text{ N}$ y trazamos a partir de su punta líneas paralelas a las direcciones de las dos sogas, estas líneas tardarán mucho en intersecarse con los ejes de acción originales. Esto genera un paralelogramo extremadamente ancho y aplanado. Para que la diagonal vertical mantenga su valor de $50\text{ N}$ mientras el ángulo interno se ensancha, los lados del paralelogramo deben estirarse enormemente. Físicamente, esto significa que la tensión necesaria para generar la misma componente vertical de elevación crece de manera desproporcionada a medida que la soga se vuelve más horizontal.

\subsubsection*{Demostración Computacional del Paralelogramo Aplanado}
El código a continuación ilustra esta situación extrema. Al trazar el paralelogramo con un ángulo muy obtuso, se evidencia de forma contundente cómo los vectores de tensión superan ampliamente la magnitud del vector resultante.

\begin{lstlisting}[language=Python]
import matplotlib.subplots as plt

fig, ax = plt.subplots(figsize=(8, 4))

ax.annotate("", xy=(0, 50), xytext=(0, 0), arrowprops=dict(facecolor='blue', width=3, headwidth=10))
ax.text(5, 25, "Resultante = 50 N", color="blue", fontsize=10, fontweight='bold')

ax.annotate("", xy=(-80, 25), xytext=(0, 0), arrowprops=dict(facecolor='red', width=2, headwidth=8))
ax.annotate("", xy=(80, 25), xytext=(0, 0), arrowprops=dict(facecolor='red', width=2, headwidth=8))
ax.text(-85, 12, "Tension >> 50 N", color="red", fontsize=10)
ax.text(50, 12, "Tension >> 50 N", color="red", fontsize=10)

ax.plot([-80, 0], [25, 50], color='gray', linestyle='--')
ax.plot([80, 0], [25, 50], color='gray', linestyle='--')

ax.set_xlim(-100, 100)
ax.set_ylim(-5, 60)
plt.title("Regla del Paralelogramo: Angulo Muy Obtuso")
plt.axis("off")
plt.show()
\end{lstlisting}

\subsubsection*{Conclusión}
La tensión en la soga es drásticamente mayor que $50\text{ N}$. La regla del paralelogramo demuestra que, al formar un ángulo muy abierto, los lados requeridos para completar la figura y alcanzar la diagonal de $50\text{ N}$ son inmensamente más largos que la diagonal misma.

\newpage

\subsection*{Enunciado}
Tu amiga se sienta en reposo sobre una silla. ¿Puedes decir que ninguna fuerza actúa sobre ella? ¿O es correcto decir que ninguna fuerza neta actúa sobre ella? Defiende tu respuesta.

\subsection*{Solución}

\subsubsection*{Diferencia entre Fuerza y Fuerza Neta}
De acuerdo con la Primera Ley de Newton y los conceptos de equilibrio estático expuestos por Paul G. Hewitt en Física Conceptual, es fundamental distinguir entre las fuerzas individuales que interactúan con un cuerpo y la fuerza resultante (o fuerza neta) de la combinación de todas ellas. Un objeto que se encuentra en estado de reposo no implica la ausencia mágica de fuerzas, sino que señala un balance perfecto entre las distintas interacciones mecánicas presentes.

\subsubsection*{Análisis de las Fuerzas en Juego}
Cuando tu amiga está sentada tranquilamente en la silla, existe una fuerza constante de largo alcance actuando sobre su centro de masa en todo momento: la atracción gravitacional de la Tierra, comúnmente denominada peso. Si esta fuera la única fuerza, ella aceleraría en caída libre hacia el centro de la Tierra. Sin embargo, los átomos que componen la superficie del asiento de la silla son comprimidos microscópicamente bajo ella y reaccionan empujando hacia arriba de forma perpendicular a la superficie con una fuerza de soporte (conocida como fuerza normal). 
Estas dos fuerzas concurrentes (el peso dirigido hacia abajo y la fuerza normal dirigida hacia arriba) son completamente iguales en magnitud pero estrictamente opuestas en dirección. Al realizar la suma vectorial de ambas, se cancelan mutuamente, arrojando como resultado algebraico un valor de cero ($\Sigma F = 0$).

\subsubsection*{Conclusión}
Es incorrecto decir que ninguna fuerza actúa sobre ella, ya que la fuerza de gravedad y la fuerza de soporte de la silla están presentes e interactuando continuamente con su cuerpo. Lo correcto y físicamente riguroso es afirmar que ninguna fuerza neta actúa sobre ella, debido a que el sistema se encuentra en equilibrio mecánico y la suma vectorial de las fuerzas individuales se anula.

\newpage

\subsection*{Enunciado}
Nellie Newton cuelga en reposo de los extremos de la soga que se muestra. ¿Cómo se compara la lectura de la báscula con su peso?

\begin{center}
\begin{tikzpicture}[scale=0.8]
  \draw[thick] (-1.5, 2.5) -- (1.5, 2.5);
  \foreach \x in {-1.4, -1.0, -0.6, -0.2, 0.2, 0.6, 1.0, 1.4}
    \draw (\x, 2.5) -- (\x+0.2, 2.7);
  
  \draw[thick] (0, 2.5) -- (0, 1.5);
  \draw[thick, fill=white] (0, 1.5) circle (1);
  \draw[thick, fill=gray!30] (-0.1, 1.5) rectangle (0.1, 2.5);
  \draw[thick, fill=white] (0, 1.5) circle (0.2);
  
  \draw[thick] (-1, 1.5) -- (-1, -1);
  
  \draw[thick] (1, 1.5) -- (1, 0.5);
  \draw[thick, fill=white, rounded corners] (0.8, -0.5) rectangle (1.2, 0.5);
  \draw[thick] (1, 0.5) -- (1, -0.5);
  \draw (0.9, 0) -- (1.1, 0);
  \draw (0.9, 0.2) -- (1.1, 0.2);
  \draw (0.9, -0.2) -- (1.1, -0.2);
  \draw[thick] (1, -0.5) -- (1, -1);
  
  \draw[thick, fill=green!40, rounded corners] (-0.8, -2.5) rectangle (0.8, -1.5);
  \draw[thick, fill=white] (0, -1.2) circle (0.4);
  \draw[thick] (-0.2, -1.1) arc (180:360:0.2);
  \draw[thick] (-0.1, -1.05) node {.};
  \draw[thick] (0.1, -1.05) node {.};
  
  \draw[thick] (-0.8, -1.5) -- (-1, -1);
  \draw[thick] (0.8, -1.5) -- (1, -1);
  
  \draw[thick, fill=cyan!50] (-0.4, -2.5) rectangle (-0.2, -2.8);
  \draw[thick, fill=cyan!50] (0.2, -2.5) rectangle (0.4, -2.8);
  \draw[thick] (-0.5, -2.8) -- (-0.1, -2.8);
  \draw[thick] (0.1, -2.8) -- (0.5, -2.8);
\end{tikzpicture}
\end{center}

\subsection*{Solución}

\subsubsection*{Análisis del Equilibrio Estático}
El enunciado estipula explícitamente que Nellie Newton se encuentra colgada "en reposo", lo que dictamina que todo el sistema carece de aceleración y reside en equilibrio mecánico estático. Si aplicamos la regla fundamental del equilibrio, sabemos que la sumatoria total de fuerzas a lo largo del eje vertical debe ser nula ($\Sigma F_{y} = 0$). En términos prácticos, esto significa que la magnitud total de la fuerza que intenta llevar a Nellie hacia abajo debe ser equilibrada por una fuerza de idéntica magnitud que tira de ella hacia arriba.

\subsubsection*{Distribución de la Tensión a través de una Polea}
La fuerza principal orientada hacia abajo es la fuerza de gravedad que actúa sobre la masa de Nellie, es decir, su peso total ($W$). En contraposición, orientada hacia arriba, no encontramos un único cable soportando toda esa carga, sino que observamos dos segmentos completamente verticales de la misma soga de los cuales ella se sostiene. 
Una polea ideal sin fricción tiene la propiedad mecánica de alterar la dirección en la que se aplica una fuerza, pero preservando intacta su magnitud. De esto se deduce que la tensión interna se distribuye de manera uniforme y constante a lo largo de toda la longitud de la soga. Por consiguiente, los dos segmentos de soga que sostienen a Nellie ejercen un tirón ascendente y cada uno contribuye con una tensión idéntica ($T$).

\subsubsection*{Diagrama de Cuerpo Libre Computacional}
Para consolidar la percepción de cómo la carga gravitacional se divide simétricamente entre los apoyos, se ofrece el siguiente código de modelado. En el diagrama, el vector singular que representa el peso descendente es contrarrestado armónicamente por dos vectores de tensión de menor envergadura.

\begin{lstlisting}[language=Python]
import matplotlib.subplots as plt

fig, ax = plt.subplots(figsize=(6, 7))

ax.plot(0, 0, 'ko', markersize=12)

ax.annotate("", xy=(0, -200), xytext=(0, 0), arrowprops=dict(facecolor='blue', width=4, headwidth=12))
ax.text(0.2, -100, "W (Peso total de Nellie)", color="blue", fontsize=11, fontweight='bold')

ax.annotate("", xy=(-0.5, 100), xytext=(0, 0), arrowprops=dict(facecolor='red', width=4, headwidth=12))
ax.text(-2.0, 50, "T (Soga Izq.)", color="red", fontsize=10)

ax.annotate("", xy=(0.5, 100), xytext=(0, 0), arrowprops=dict(facecolor='red', width=4, headwidth=12))
ax.text(0.7, 50, "T (Soga Der.)\nLectura de Bascula", color="red", fontsize=10)

ax.set_xlim(-2.5, 2.5)
ax.set_ylim(-250, 150)
plt.title("DCL: Equilibrio de Tensiones Paralelas")
plt.axis("off")
plt.show()
\end{lstlisting}

\subsubsection*{Planteamiento Matemático}
Para encontrar la lectura, estructuramos la ecuación de equilibrio sumando las fuerzas ascendentes y la igualamos a la fuerza descendente:
$$T + T = W$$
Agrupando los términos equivalentes:
$$2T = W$$
Despejando la incógnita de la tensión presente en un solo segmento de soga ($T$):
$$T = \frac{W}{2}$$
Finalmente, inspeccionando la disposición del sistema mecánico en la imagen, podemos constatar que el instrumento de medición (la báscula de resorte) se encuentra insertado de manera secuencial a lo largo de la trayectoria del segmento derecho de la soga. Las básculas de este tipo miden únicamente la deformación interna causada por la fuerza de tensión que transita a través de ellas.

\subsubsection*{Conclusión}
La lectura de la báscula registrará exactamente la mitad del peso total de Nellie Newton.

\newpage

\subsection*{Enunciado}
Harry el pintor se balancea año tras año en su guindola. Su peso es de $500\text{ N}$ y la soga tiene un punto de rotura de $300\text{ N}$, que él desconoce. ¿Por qué la soga no se rompe cuando se sostiene como se muestra a la izquierda? Un día, Harry pinta cerca de un asta de bandera y, para cambiar, amarra el extremo libre de la soga al asta en lugar de a su silla, como se muestra a la derecha. ¿Por qué Harry termina tomando sus vacaciones más pronto?

\begin{center}
\begin{tikzpicture}[scale=0.8]
  \draw[thick, fill=gray!30] (0, -1) rectangle (1, 6);
  \draw[thick] (0, 5) -- (-1, 5) -- (-1, 4.8);
  \draw[thick, fill=white] (-1, 4.4) circle (0.4);
  \draw[thick] (-0.6, 4.4) -- (-0.6, 0.5);
  \draw[thick] (-1.4, 4.4) -- (-1.4, 1.2);
  \draw[thick, fill=blue!20] (-1.8, 0) rectangle (-0.2, 0.5);
  \draw[thick] (-1, 0.5) circle (0.3);
  \draw[thick] (-1, 0.2) -- (-1, -0.2);
  \draw[thick] (-1.4, 1.2) -- (-0.8, 0.8);
  \node at (-1, -1.5) {Situación Izquierda};

  \begin{scope}[xshift=7cm]
    \draw[thick, fill=gray!30] (0, -1) rectangle (1, 6);
    \draw[thick] (0, 5) -- (-1, 5) -- (-1, 4.8);
    \draw[thick, fill=white] (-1, 4.4) circle (0.4);
    \draw[thick] (-0.6, 4.4) -- (-0.6, 0.5);
    \draw[thick] (0, 1.5) -- (-4, 1.5);
    \draw[thick] (-1.4, 4.4) -- (-1.4, 1.5);
    \draw[thick, fill=blue!20] (-1.8, 0) rectangle (-0.2, 0.5);
    \draw[thick] (-1, 0.5) circle (0.3);
    \draw[thick] (-1, 0.2) -- (-1, -0.2);
    \draw[thick, fill=red!60] (-4, 1.5) rectangle (-2.5, 0);
    \draw[thick, fill=white] (-3.8, 1.3) rectangle (-3.5, 1);
    \node at (-1.5, -1.5) {Situación Derecha};
  \end{scope}
\end{tikzpicture}
\end{center}

\subsection*{Solución}

\subsubsection*{El Concepto del Límite del Sistema y la Tensión}
Para entender la mecánica de poleas y cuerdas expuesta en la Física Conceptual de Hewitt, es imperativo definir claramente el "sistema" físico que estamos analizando. Nuestro sistema está compuesto exclusivamente por Harry y su guindola (silla), cuyo peso total combinado ejerce una fuerza descendente de $500\text{ N}$. 
El equilibrio estático ($\Sigma F_{y} = 0$) exige que la fuerza ascendente total sobre este sistema sea de exactamente $500\text{ N}$. La clave del problema radica en contar cuántos segmentos de cuerda atraviesan el límite de nuestro sistema tirando hacia arriba. A través de una polea ideal (sin fricción), la tensión es uniforme en toda la cuerda.

\subsubsection*{Análisis de la Situación Izquierda}
En el primer escenario, Harry sujeta un extremo de la cuerda mientras el otro está atado a su silla. Si dibujamos un círculo imaginario alrededor de Harry y la silla, notaremos que dos segmentos verticales de la cuerda atraviesan este límite hacia arriba. Ambos segmentos están soportando el peso del sistema simultáneamente.
La ecuación de equilibrio se plantea así:
$$T_{1} + T_{2} - W = 0$$
Como la tensión es la misma en toda la cuerda ($T_{1} = T_{2} = T$):
$$2T = 500\text{ N}$$
$$T = \frac{500\text{ N}}{2}$$
$$T = 250\text{ N}$$
La tensión a la que está sometida la cuerda es de $250\text{ N}$. Dado que el límite de rotura de la soga es de $300\text{ N}$, y $250\text{ N} < 300\text{ N}$, la cuerda soporta el esfuerzo sin romperse.

\subsubsection*{Análisis de la Situación Derecha}
En el segundo escenario, Harry amarra el extremo libre de la cuerda al asta de la bandera. Si volvemos a dibujar el límite imaginario alrededor de Harry y su silla, observamos un cambio drástico: ahora solo un segmento de la cuerda (el que está atado a la silla) atraviesa el límite tirando hacia arriba. El otro segmento está soportando el asta de la bandera, una fuerza externa que no contribuye a elevar a Harry.
La nueva ecuación de equilibrio para el sistema de Harry es:
$$T - W = 0$$
$$T = 500\text{ N}$$
En este caso, un solo segmento debe soportar toda la carga. La tensión resultante es de $500\text{ N}$. Como este valor excede ampliamente el límite de rotura de $300\text{ N}$ ($500\text{ N} > 300\text{ N}$), la estructura interna de la cuerda colapsa.

\subsubsection*{Diagrama de Cuerpo Libre Computacional}
El siguiente script en Python genera dos diagramas que contrastan ambas situaciones, ilustrando matemáticamente por qué la distribución de fuerzas define el fracaso estructural en el segundo caso.

\begin{lstlisting}[language=Python]
import matplotlib.pyplot as plt

fig, (ax1, ax2) = plt.subplots(1, 2, figsize=(10, 6))

ax1.plot(0, 0, 'ko', markersize=15)
ax1.annotate("", xy=(0, -500), xytext=(0, 0), arrowprops=dict(facecolor='blue', width=4))
ax1.text(0.2, -250, "W = 500 N", color="blue", fontsize=12)
ax1.annotate("", xy=(-0.5, 250), xytext=(0, 0), arrowprops=dict(facecolor='green', width=4))
ax1.text(-2, 125, "T = 250 N", color="green", fontsize=12)
ax1.annotate("", xy=(0.5, 250), xytext=(0, 0), arrowprops=dict(facecolor='green', width=4))
ax1.text(0.7, 125, "T = 250 N", color="green", fontsize=12)
ax1.set_title("Izquierda: 2 Segmentos (Seguro)")
ax1.set_xlim(-3, 3)
ax1.set_ylim(-600, 600)
ax1.axis("off")

ax2.plot(0, 0, 'ko', markersize=15)
ax2.annotate("", xy=(0, -500), xytext=(0, 0), arrowprops=dict(facecolor='blue', width=4))
ax2.text(0.2, -250, "W = 500 N", color="blue", fontsize=12)
ax2.annotate("", xy=(0, 500), xytext=(0, 0), arrowprops=dict(facecolor='red', width=4))
ax2.text(0.2, 250, "T = 500 N\n(¡Excede 300 N!)", color="red", fontsize=12, fontweight='bold')
ax2.set_title("Derecha: 1 Segmento (Rotura)")
ax2.set_xlim(-3, 3)
ax2.set_ylim(-600, 600)
ax2.axis("off")

plt.tight_layout()
plt.show()
\end{lstlisting}

\subsubsection*{Conclusión}
A la izquierda, la soga no se rompe porque la fuerza gravitacional de $500\text{ N}$ se distribuye equitativamente entre los dos segmentos verticales de la cuerda, generando una tensión de $250\text{ N}$ (inferior al límite de rotura). Harry termina tomando vacaciones pronto a la derecha porque, al amarrar la cuerda al asta, su peso recae sobre un único segmento vertical, elevando la tensión a $500\text{ N}$, lo que supera la tolerancia de $300\text{ N}$ de la soga y provoca su ruptura.

\newpage

\subsection*{Enunciado}
Para el sistema de poleas que se muestra, ¿cuál es el límite superior de peso que puede levantar el hombre fuerte?

\begin{center}
\begin{tikzpicture}[scale=0.8]
  \draw[thick] (-3, 4) -- (3, 4);
  \foreach \x in {-2.8, -2.4, -2.0, -1.6, -1.2, -0.8, -0.4, 0, 0.4, 0.8, 1.2, 1.6, 2.0, 2.4, 2.8}
    \draw (\x, 4) -- (\x+0.2, 4.2);
    
  \draw[thick, fill=gray!30] (-0.2, 4) rectangle (0.2, 2.5);
  \draw[thick, fill=cyan!40] (0, 2.5) circle (1.2);
  \draw[thick, fill=white] (0, 2.5) circle (0.15);
  
  \draw[thick] (-1.2, 2.5) -- (-1.2, -1);
  \draw[thick, fill=cyan!20] (-2.2, -2.5) rectangle (-0.2, -1);
  \node[font=\Large\bfseries] at (-1.2, -1.75) {W};
  \draw[thick] (-1.2, -1) arc (-90:90:0.15);
  
  \draw[thick] (1.2, 2.5) -- (1.2, 0.5);
  
  \draw[thick, fill=white] (1.5, 0.2) circle (0.3); 
  \draw[thick] (1.5, -0.1) -- (1.5, -1.2);
  \draw[thick] (1.5, -1.2) -- (1.0, -2.5);
  \draw[thick] (1.5, -1.2) -- (2.0, -2.5);
  \draw[thick] (1.5, -0.3) -- (1.2, 0.5);
  \draw[thick] (1.5, -0.5) -- (1.2, 0.3);
  
  \draw[thick] (-3, -2.5) -- (3, -2.5);
\end{tikzpicture}
\end{center}

\subsection*{Solución}

\subsubsection*{Mecánica de una Polea Fija Simple}
En la Física Conceptual, una polea fija simple (como la que se muestra anclada al techo) funciona estrictamente como un cambiador de dirección. No multiplica la fuerza; su ventaja mecánica teórica es igual a 1. Esto significa que la tensión que se aplica en un extremo de la cuerda se transmite íntegramente a lo largo de la misma, ejerciendo exactamente la misma fuerza de tracción en el extremo opuesto, pero en sentido contrario. Para levantar el bloque $W$, el hombre debe tirar de la soga hacia abajo con una fuerza igual al peso del bloque.

\subsubsection*{Tercera Ley de Newton y el Límite Físico}
Para determinar el límite de fuerza que el hombre puede aplicar, debemos analizar las interacciones de fuerza sobre él mismo utilizando la Tercera Ley de Newton (acción y reacción). Cuando el hombre tira de la soga hacia abajo con una cierta fuerza (tensión $T$), la soga reacciona tirando del hombre hacia arriba con exactamente esa misma fuerza $T$. 
Al mismo tiempo, la gravedad tira del hombre hacia abajo con una fuerza igual a su propio peso ($W_{\text{hombre}}$). Para que el hombre permanezca en el suelo y pueda seguir tirando, la fuerza neta hacia abajo debe ser mayor o igual a cero. Si la soga tira de él hacia arriba con una fuerza superior a su propio peso ($T > W_{\text{hombre}}$), la fuerza normal del piso desaparecerá y el hombre será levantado del suelo. 

\subsubsection*{Diagrama de Cuerpo Libre Computacional}
El siguiente script modela el equilibrio límite del hombre. Muestra visualmente que el vector de tensión hacia arriba no puede exceder el vector de su peso hacia abajo, de lo contrario, el hombre aceleraría hacia arriba en lugar del bloque.

\begin{lstlisting}[language=Python]
import matplotlib.pyplot as plt

fig, ax = plt.subplots(figsize=(5, 7))

ax.plot(0, 0, 'ko', markersize=12)

ax.annotate("", xy=(0, -800), xytext=(0, 0), arrowprops=dict(facecolor='blue', width=4, headwidth=12))
ax.text(0.2, -400, "W_hombre", color="blue", fontsize=12, fontweight='bold')

ax.annotate("", xy=(0, 800), xytext=(0, 0), arrowprops=dict(facecolor='red', width=4, headwidth=12))
ax.text(0.2, 400, "T_max (Limite) = W_hombre", color="red", fontsize=11)

ax.set_xlim(-2, 2)
ax.set_ylim(-1000, 1000)
plt.title("DCL del Hombre Fuerte: Caso Limite")
plt.axis("off")
plt.show()
\end{lstlisting}

\subsubsection*{Conclusión}
El límite superior de peso que puede levantar el hombre fuerte es su propio peso. Si intenta levantar un objeto más pesado, la fuerza necesaria superará su peso corporal y terminará izándose a sí mismo hacia el techo en lugar de levantar el bloque.

\newpage

\subsection*{Enunciado}
Si el hombre fuerte del ejercicio anterior ejerce una fuerza descendente de $800\text{ N}$ sobre la soga, ¿cuánta fuerza ascendente se ejerce sobre el cubo?

\subsection*{Solución}

\subsubsection*{Principio de Transmisión de Tensión}
Tal como se estableció en el análisis del sistema físico previo, una cuerda ideal (asumiendo que su propia masa es despreciable y no tiene elasticidad significativa) sometida a tracción transmite la fuerza de manera idéntica y sin pérdidas a través de toda su longitud. La polea fija, al carecer de fricción en este modelo conceptual, se limita exclusivamente a redirigir esta fuerza un ángulo de 180 grados sin alterar su magnitud escalar.

\subsubsection*{Relación Directa de Fuerzas}
El hombre interactúa con su extremo de la soga ejerciendo una fuerza de $800\text{ N}$ en dirección descendente. Por la Tercera Ley de Newton, la soga adquiere una tensión interna ($T$) de exactamente $800\text{ N}$. 
A medida que esta tensión se propaga por el segmento de la cuerda, pasa sobre la rueda de la polea y desciende hacia el bloque, mantiene su valor inalterado. En el punto de anclaje con el cubo $W$, esta misma cuerda está sometida a la misma tensión de $800\text{ N}$, pero debido a la geometría impuesta por la polea, su dirección ahora es puramente ascendente.

\subsubsection*{Representación Computacional de la Fuerza sobre el Cubo}
Para evitar ambigüedades, el diagrama de cuerpo libre del cubo confirma que la única fuerza de soporte presente es generada por el segmento único de cuerda que lo ata a la polea.

\begin{lstlisting}[language=Python]
import matplotlib.pyplot as plt

fig, ax = plt.subplots(figsize=(4, 5))

ax.plot(0, 0, 'ks', markersize=20)

ax.annotate("", xy=(0, 800), xytext=(0, 0), arrowprops=dict(facecolor='magenta', width=4, headwidth=12))
ax.text(0.2, 400, "T = 800 N", color="magenta", fontsize=12, fontweight='bold')

ax.set_xlim(-2, 2)
ax.set_ylim(-200, 1000)
plt.title("DCL: Fuerza Ascendente sobre el Cubo")
plt.axis("off")
plt.show()
\end{lstlisting}

\subsubsection*{Conclusión}
Se ejerce una fuerza ascendente de $800\text{ N}$ sobre el cubo.

\newpage

\subsection*{Enunciado}
¿Cuántas fuerzas significativas actúan sobre un libro en reposo en una mesa? Identifica las fuerzas.

\subsection*{Solución}

\subsubsection*{El Concepto de Equilibrio Estático}
De acuerdo con la Física Conceptual de Paul G. Hewitt, cuando un objeto se encuentra en reposo absoluto sobre una superficie, se dice que está en un estado de equilibrio estático. En este estado, la Primera Ley de Newton dicta que la sumatoria de todas las fuerzas que actúan sobre el cuerpo debe ser igual a cero ($\Sigma F = 0$). Para que esto ocurra, las fuerzas presentes deben equilibrarse mutuamente.

\subsubsection*{Identificación de las Interacciones Físicas}
A nivel macroscópico, existen dos interacciones fundamentales sobre el libro:
1. La primera es una fuerza de campo de largo alcance: la atracción gravitacional que la masa de la Tierra ejerce sobre la masa del libro. Esta fuerza empuja invariablemente al libro hacia el centro del planeta, y es lo que comúnmente llamamos el peso del libro.
2. La segunda es una fuerza de contacto mecánico. Dado que el libro no atraviesa la materia sólida de la mesa, los átomos de la superficie de la mesa repelen eléctricamente a los átomos de la cubierta del libro. Esta fuerza de reacción estructural empuja al libro de forma perpendicular a la superficie, alejándolo de ella, y se conoce como fuerza normal o fuerza de soporte.

\subsubsection*{Diagrama de Cuerpo Libre Computacional}
El siguiente código genera una representación visual del libro como una partícula puntual, demostrando cómo ambas fuerzas actúan a lo largo del mismo eje vertical pero con sentidos diametralmente opuestos para lograr el equilibrio.

\begin{lstlisting}[language=Python]
import matplotlib.pyplot as plt

fig, ax = plt.subplots(figsize=(4, 5))

ax.plot(0, 0, 'ks', markersize=40)

ax.annotate("", xy=(0, -50), xytext=(0, 0), arrowprops=dict(facecolor='blue', width=4))
ax.text(0.15, -25, "W (Gravedad)", color="blue", fontsize=12, fontweight='bold')

ax.annotate("", xy=(0, 50), xytext=(0, 0), arrowprops=dict(facecolor='red', width=4))
ax.text(0.15, 25, "N (Soporte)", color="red", fontsize=12, fontweight='bold')

ax.set_xlim(-1, 1)
ax.set_ylim(-80, 80)
plt.title("Fuerzas sobre un libro en reposo")
plt.axis("off")
plt.show()
\end{lstlisting}

\subsubsection*{Conclusión}
Actúan dos fuerzas significativas: la fuerza de gravedad (que tira del libro hacia abajo) y la fuerza de soporte o normal (que la mesa ejerce empujando el libro hacia arriba).

\newpage
\subsection*{Enunciado}
Coloca un libro pesado sobre una mesa y la mesa empuja hacia arriba sobre el libro. ¿Por qué este empuje hacia arriba no hace que el libro se eleve de la mesa?

\subsection*{Solución}

\subsubsection*{Relación entre Fuerza Neta y Aceleración}
La Segunda Ley de Newton establece que la aceleración de un objeto es directamente proporcional a la fuerza neta que actúa sobre él. Para que el libro adquiera velocidad hacia arriba (es decir, para que se eleve desde el reposo), debe existir una fuerza neta ascendente mayor que cero.

\subsubsection*{El Principio de Cancelación de Fuerzas}
Como analizamos previamente, la fuerza de empuje de la mesa hacia arriba (fuerza normal) es una respuesta pasiva a la compresión estructural que sufre la mesa bajo el peso del libro. La mesa solo empuja hacia arriba con la fuerza exacta que se necesita para contrarrestar la fuerza descendente de la gravedad. 
Matemáticamente, si asignamos valores, la fuerza hacia arriba es exactamente igual en magnitud a la fuerza hacia abajo. Al sumarlas vectorialmente, el resultado aritmético es cero. 
Al no haber una fuerza neta o "sobrante" hacia arriba, la aceleración vertical es matemáticamente cero, imposibilitando que el objeto altere su estado de reposo y despegue de la superficie.

\subsubsection*{Conclusión}
El empuje hacia arriba no eleva el libro porque dicha fuerza está perfectamente equilibrada y anulada por la fuerza de gravedad descendente. Al cancelarse mutuamente, la fuerza neta es cero, por lo que no hay aceleración ascendente.

\newpage

\subsection*{Enunciado}
Cuando estás de pie sobre el piso, ¿el piso ejerce una fuerza hacia arriba contra tus pies? ¿Cuánta fuerza ejerce? ¿Por qué no te mueves hacia arriba por esta fuerza?

\subsection*{Solución}

\subsubsection*{Tercera Ley de Newton (Acción y Reacción)}
Según la Tercera Ley del movimiento, las fuerzas siempre ocurren en pares. La fuerza no es algo que un objeto posee, sino una interacción entre dos objetos. Cuando te paras sobre el piso, interactúas con él. Tu masa, arrastrada por la gravedad, ejerce una fuerza hacia abajo contra los azulejos del piso (la acción). Instantánea y simultáneamente, el piso responde ejerciendo una fuerza hacia arriba contra las plantas de tus pies (la reacción).

\subsubsection*{Cuantificación de la Fuerza de Reacción}
Asumiendo que estás de pie en completo reposo estático, se debe cumplir la regla del equilibrio ($\Sigma F = 0$). Para que tu cuerpo no se hunda en el suelo ni salga proyectado hacia el techo, la suma de las fuerzas verticales debe anularse. Por ende, la fuerza normal que proporciona la estructura de concreto del piso debe ser matemáticamente equivalente a la fuerza con la que la Tierra tira de ti hacia abajo.

\subsubsection*{Dinámica de la Ausencia de Movimiento}
La razón por la que no te mueves hacia arriba a pesar de este constante y poderoso empuje ascendente se debe al concepto de fuerza neta abordado en el problema anterior. La fuerza del piso es solo la mitad de la historia. Todo el tiempo que el piso te empuja hacia arriba, la gravedad te empuja hacia abajo con igual intensidad. Ambas fuerzas actúan sobre el mismo cuerpo (el tuyo), resultando en una red de fuerzas equivalente a cero.

\subsubsection*{Conclusión}
Sí, el piso ejerce una fuerza hacia arriba. Ejerce una cantidad de fuerza exactamente igual a la de tu peso. No te mueves hacia arriba porque esta fuerza ascendente está completamente contrarrestada y anulada por tu peso dirigido hacia abajo, manteniendo tu fuerza neta en cero.

\newpage

\subsection*{Enunciado}
Supón que rebotas arriba y abajo mientras te pesas en una báscula de baño. ¿Qué varía: la fuerza de sostén ascendente o la fuerza de gravedad sobre ti? ¿Por qué la lectura de tu peso se muestra mejor cuando estás de pie en reposo sobre la báscula?

\subsection*{Solución}

\subsubsection*{Análisis del Componente Constante: La Gravedad}
La fuerza de gravedad (el peso real) se define por la ecuación $W = m \cdot g$. Tu masa corporal ($m$) no cambia por el simple hecho de saltar. Asimismo, la aceleración debida a la gravedad terrestre ($g \approx 9.8 \text{ m/s}^2$) es una constante del campo planetario. Puesto que ninguna de las dos variables matemáticas que definen la gravedad se altera, la fuerza que la Tierra ejerce para tirar de ti hacia su centro permanece estrictamente constante, sin importar lo que hagas encima de la báscula.

\subsubsection*{Análisis del Componente Variable: La Fuerza de Sostén}
Cuando rebotas, cambias la velocidad de tu centro de masa; en otras palabras, estás acelerando hacia arriba y hacia abajo alternadamente. 
Para acelerar hacia arriba desde un punto bajo, necesitas que la fuerza neta apunte hacia arriba ($\Sigma F > 0$). Dado que la gravedad es constante, la única forma de que la red de fuerzas sea positiva hacia arriba es que la báscula te empuje con una fuerza de sostén mayor que tu propio peso. 
Por el contrario, para acelerar hacia abajo, la fuerza neta debe apuntar hacia el suelo. Esto requiere que la báscula disminuya su presión contra tus pies, proporcionando una fuerza de sostén menor que tu peso. 

\subsubsection*{Mecánica de Medición de una Báscula}
Es un error común pensar que las básculas "leen" la gravedad. Como se explica en la Física Conceptual, una báscula de baño no es un medidor de gravedad, es un medidor de deformación mecánica. Mide estrictamente la fuerza de compresión a la que está sometida su superficie superior, que no es otra cosa que la fuerza de sostén (fuerza normal) que ella misma debe ejercer.

\subsubsection*{Representación Computacional de las Fuerzas Dinámicas}
El siguiente script modela el escenario de una persona rebotando sobre una báscula. Grafica cómo la fuerza constante de la gravedad se mantiene como una línea recta, mientras que la lectura de la báscula (la fuerza de sostén) oscila violentamente debido a la aceleración del rebote.

\begin{lstlisting}[language=Python]
import matplotlib.pyplot as plt
import numpy as np

t = np.linspace(0, 10, 200)
W_real = np.full_like(t, 700) 
N_bascula = 700 + 300 * np.sin(2 * np.pi * 0.5 * t) 

fig, ax = plt.subplots(figsize=(8, 4))

ax.plot(t, W_real, 'b-', label='Gravedad (Peso Real = 700 N)', linewidth=3)
ax.plot(t, N_bascula, 'r--', label='Lectura de Bascula (Fuerza de Sosten)', linewidth=2)

ax.set_xlabel("Tiempo (s)")
ax.set_ylabel("Fuerza (N)")
ax.set_title("Dinamica de Fuerzas: Rebotando en una Bascula")
ax.legend(loc='lower left')
ax.grid(True, linestyle=':', alpha=0.7)

plt.show()
\end{lstlisting}

\subsubsection*{Conclusión}
Lo que varía violentamente es la fuerza de sostén ascendente; la fuerza de gravedad permanece invariable. La lectura de tu peso es exacta solo cuando estás en reposo porque el reposo garantiza que la aceleración sea cero, obligando matemáticamente a la fuerza de sostén (que es lo que la báscula realmente mide) a igualar exactamente a tu fuerza de gravedad constante.

\newpage
\subsection*{Enunciado}
Un jarro vacío con peso $W$ reposa sobre una mesa. ¿Cuál es la fuerza de sostén que ejerce la mesa sobre el jarro?

\subsection*{Solución}

\subsubsection*{Condición de Equilibrio Estático}
Según la Primera Ley de Newton tratada en la Física Conceptual, cuando un objeto se encuentra en reposo absoluto sobre una superficie, su estado cinemático indica que no posee aceleración. En consecuencia, la fuerza neta o resultante que actúa sobre él es cero ($\Sigma F = 0$).

\subsubsection*{Análisis de Interacciones}
En este escenario unidimensional en el eje vertical, el jarro experimenta dos fuerzas. Hacia abajo actúa la fuerza de atracción gravitacional de la Tierra, que el enunciado define como el peso $W$. Para que el jarro no acelere hacia abajo y atraviese la mesa, esta última debe reaccionar empujando hacia arriba con una fuerza de igual magnitud. Esta fuerza de reacción mecánica de la superficie es la fuerza de sostén o fuerza normal.

\subsubsection*{Planteamiento Matemático}
Aplicando la sumatoria de fuerzas en el eje vertical:
$$\Sigma F_{y} = 0$$
$$N - W = 0$$
$$N = W$$

\subsubsection*{Conclusión}
La fuerza de sostén que ejerce la mesa sobre el jarro es exactamente igual a $W$.

\newpage

\subsection*{Enunciado}
Un jarro vacío con peso $W$ reposa sobre una mesa. ¿Cuál es la fuerza de sostén cuando se vierte agua con peso $w$ en el jarro?

\subsection*{Solución}

\subsubsection*{Redefinición del Sistema Físico}
Al añadir agua al jarro, debemos expandir nuestro sistema de análisis. Ahora, la carga total que la mesa debe soportar no es solo la masa del cristal o cerámica del jarro, sino también la masa del fluido que contiene. El peso es una propiedad aditiva, por lo que la fuerza total de gravedad dirigida hacia abajo será la suma de los pesos individuales de los componentes del sistema.

\subsubsection*{Diagrama de Cuerpo Libre Computacional}
El siguiente código modela las fuerzas que actúan sobre la superficie de la mesa. Muestra cómo la fuerza normal debe equilibrar la suma vectorial de los pesos del jarro y del agua de forma simultánea.

\begin{lstlisting}[language=Python]
import matplotlib.pyplot as plt

fig, ax = plt.subplots(figsize=(4, 6))

ax.plot(0, 0, 'ko', markersize=15)

ax.annotate("", xy=(0, -2), xytext=(0, 0), arrowprops=dict(facecolor='blue', width=4))
ax.text(0.15, -1, "W (Jarro)", color="blue", fontsize=11)

ax.annotate("", xy=(0, -4), xytext=(0, -2), arrowprops=dict(facecolor='cyan', width=4))
ax.text(0.15, -3.5, "w (Agua)", color="cyan", fontsize=11)

ax.annotate("", xy=(0, 4), xytext=(0, 0), arrowprops=dict(facecolor='red', width=4))
ax.text(0.15, 2, "N = W + w", color="red", fontsize=11, fontweight='bold')

ax.set_xlim(-2, 2)
ax.set_ylim(-5, 5)
plt.title("DCL: Jarro conteniendo Agua")
plt.axis("off")
plt.show()
\end{lstlisting}

\subsubsection*{Ecuación de Equilibrio Ajustada}
La sumatoria de fuerzas verticales se modifica para incluir ambos pesos apuntando hacia abajo:
$$\Sigma F_{y} = 0$$
$$N - W - w = 0$$
Despejando la fuerza normal ($N$):
$$N = W + w$$

\subsubsection*{Conclusión}
La fuerza de sostén cuando se vierte el agua es igual a la suma de los pesos, es decir, $W + w$.

\newpage

\subsection*{Enunciado}
Si jalas horizontalmente una caja con una fuerza de $200\text{ N}$, ésta se desliza por el piso con equilibrio dinámico. ¿Cuánta fricción actúa sobre la caja?

\subsection*{Solución}

\subsubsection*{El Concepto de Equilibrio Dinámico}
Uno de los conceptos más importantes en la Física Conceptual de Hewitt es que el equilibrio no es exclusivo del reposo. El equilibrio mecánico ocurre siempre que no hay cambios en el estado de movimiento, lo cual significa que la aceleración es cero. Cuando un objeto se desliza a velocidad constante en una línea recta (sin acelerar ni frenar), se dice que está en "equilibrio dinámico". La regla fundamental sigue siendo la misma: la fuerza neta debe ser cero ($\Sigma F = 0$).

\subsubsection*{Interacción de Fuerzas Horizontales}
A lo largo del eje horizontal de movimiento, existen dos fuerzas opuestas interactuando. Hacia adelante actúa la fuerza de tensión o jalón aplicada activamente, que es de $200\text{ N}$. Hacia atrás actúa la fuerza de fricción cinética, generada por las irregularidades microscópicas entre la base de la caja y la superficie del piso. Para que la caja no acelere a pesar de estar siendo jalada, la fricción debe neutralizar exactamente el jalón.

\subsubsection*{Diagrama de Cuerpo Libre Computacional}
El siguiente script evidencia cómo las fuerzas en el eje horizontal se cancelan mutuamente para garantizar que no exista una fuerza neta que provoque aceleración.

\begin{lstlisting}[language=Python]
import matplotlib.pyplot as plt

fig, ax = plt.subplots(figsize=(6, 4))

ax.plot(0, 0, 'ks', markersize=35)

ax.annotate("", xy=(200, 0), xytext=(0, 0), arrowprops=dict(facecolor='green', width=4, headwidth=12))
ax.text(60, 25, "F_aplicada = 200 N", color="green", fontsize=11, fontweight='bold')

ax.annotate("", xy=(-200, 0), xytext=(0, 0), arrowprops=dict(facecolor='orange', width=4, headwidth=12))
ax.text(-220, 25, "f_friccion = ?", color="orange", fontsize=11, fontweight='bold')

ax.set_xlim(-300, 300)
ax.set_ylim(-100, 100)
plt.title("DCL: Equilibrio Dinamico (Eje X)")
plt.axis("off")
plt.show()
\end{lstlisting}

\subsubsection*{Resolución Algebraica}
Planteamos la condición de equilibrio para el eje x:
$$\Sigma F_{x} = 0$$
$$F_{\text{aplicada}} - f_{\text{fricción}} = 0$$
Sustituyendo el valor conocido:
$$200\text{ N} - f_{\text{fricción}} = 0$$
$$f_{\text{fricción}} = 200\text{ N}$$

\subsubsection*{Conclusión}
Actúa una fuerza de fricción de $200\text{ N}$ en sentido opuesto al movimiento.

\newpage

\subsection*{Enunciado}
Con la finalidad de deslizar sobre el piso un pesado armario con rapidez constante, ejerces una fuerza horizontal de $600\text{ N}$. ¿La fuerza de fricción entre el armario y el piso es mayor que, menor que o igual a $600\text{ N}$? Defiende tu respuesta.

\subsection*{Solución}

\subsubsection*{Análisis Cinemático del Enunciado}
La clave para resolver este problema reside en la frase "con rapidez constante". De acuerdo con la Segunda Ley de Newton ($F = m \cdot a$), la presencia de una fuerza neta desequilibrada provoca obligatoriamente una aceleración. Si la rapidez es constante y la trayectoria es recta, la aceleración es estrictamente cero. Esto implica que la sumatoria de todas las fuerzas en la dirección del movimiento debe anularse perfectamente.

\subsubsection*{Balance de Fuerzas de Contacto}
Al intentar empujar el armario, la fuerza de aplicación muscular hacia adelante es de $600\text{ N}$. La fricción cinética es una fuerza reactiva que se opone al deslizamiento de las superficies.
Si la fricción fuera menor a $600\text{ N}$, quedaría una fuerza neta hacia adelante y el armario aceleraría progresivamente, aumentando su rapidez.
Si la fricción fuera mayor a $600\text{ N}$, habría una fuerza neta hacia atrás, lo que desaceleraría el armario hasta detenerlo.
Por lo tanto, la única condición física que permite mantener la rapidez constante prometida por el enunciado es que ambas magnitudes sean idénticas.

\subsubsection*{Conclusión}
La fuerza de fricción es igual a $600\text{ N}$. Al moverse con rapidez constante, el armario se encuentra en equilibrio dinámico (aceleración cero), lo que exige que la fuerza de fricción anule exactamente la fuerza aplicada, haciendo que la fuerza neta sea cero.

\newpage

\subsection*{Enunciado}
Piensa en una caja en reposo sobre el piso de una fábrica. Cuando un par de trabajadores comienzan a levantarla, ¿la fuerza de sostén sobre la caja proporcionada por el piso aumenta, disminuye o permanece igual? ¿Qué sucede con la fuerza de sostén sobre los pies de los trabajadores?

\subsection*{Solución}

\subsubsection*{Análisis del Equilibrio en la Caja}
De acuerdo con los principios de la Física Conceptual, cuando la caja está sola y en reposo sobre el piso, la fuerza de gravedad (peso de la caja) que tira hacia abajo está perfectamente equilibrada por la fuerza de sostén (fuerza normal) que el piso ejerce hacia arriba. 
Cuando los trabajadores comienzan a aplicar una fuerza muscular hacia arriba para intentar levantarla, introducen una nueva fuerza ascendente en el sistema. Dado que la caja aún no se despega del piso, se mantiene el equilibrio estático ($\Sigma F_{y} = 0$). La ecuación de equilibrio se transforma en: Fuerza Normal del piso + Fuerza de elevación de los trabajadores = Peso de la caja. 
Al despejar la Fuerza Normal, resulta evidente que el piso ahora necesita ejercer menos fuerza para sostener la caja, ya que los trabajadores están compensando parte del peso.

\subsubsection*{La Tercera Ley de Newton en los Trabajadores}
Para comprender qué sucede con los pies de los trabajadores, debemos aplicar la Tercera Ley de Newton (acción y reacción). Para tirar de la caja hacia arriba, los trabajadores deben afianzar sus cuerpos y presionar físicamente contra el piso hacia abajo. Además, debido a la interacción, si los trabajadores tiran de la caja hacia arriba con una cierta fuerza, la caja reacciona tirando de los brazos de los trabajadores hacia abajo con esa misma cantidad de fuerza.
Por lo tanto, la fuerza total que empuja a los trabajadores hacia abajo contra el piso ahora es la suma de su propio peso corporal más la fuerza de reacción de la caja.

\subsubsection*{Visualización de la Transferencia de Cargas}
El siguiente código modela la transferencia de la fuerza de soporte. Se demuestra cómo la disminución de la fuerza normal sobre la caja es geométricamente proporcional al aumento de la fuerza normal sobre los trabajadores.

\begin{lstlisting}[language=Python]
import matplotlib.pyplot as plt

fig, axs = plt.subplots(1, 2, figsize=(10, 5))

axs[0].bar(["Caja", "Trabajadores"], [1000, 1600], color=['blue', 'green'])
axs[0].set_title("Antes de levantar (Fuerza Normal)")
axs[0].set_ylabel("Fuerza de Sosten (N)")
axs[0].set_ylim(0, 2200)

axs[1].bar(["Caja", "Trabajadores"], [600, 2000], color=['cyan', 'lime'])
axs[1].set_title("Al intentar levantar (Fuerza Normal)")
axs[1].set_ylim(0, 2200)

plt.suptitle("Transferencia de Carga en el Piso")
plt.show()
\end{lstlisting}

\subsubsection*{Conclusión}
La fuerza de sostén sobre la caja proporcionada por el piso disminuye. Por el contrario, la fuerza de sostén sobre los pies de los trabajadores aumenta.

\newpage

\subsection*{Enunciado}
Dos personas jalan cada una con una fuerza de $300\text{ N}$ sobre una soga en un concurso de tirar la cuerda. ¿Cuál es la fuerza neta sobre la soga? ¿Cuánta fuerza ejerce la soga sobre cada persona?

\subsection*{Solución}

\subsubsection*{Cálculo de la Fuerza Neta sobre el Sistema}
El concepto de fuerza neta se refiere a la suma vectorial de todas las fuerzas externas que actúan sobre un mismo objeto. En este caso, el objeto de estudio es la soga. Suponiendo que las personas están empatadas y la soga no se acelera hacia ningún lado, el sistema se encuentra en equilibrio estático. Una persona tira hacia la derecha con $300\text{ N}$ y la otra tira hacia la izquierda con $300\text{ N}$. Al sumar vectores de igual magnitud y dirección opuesta, el resultado algebraico es una cancelación total.

\subsubsection*{La Naturaleza de la Tensión y la Reacción}
Para responder a la segunda pregunta, abandonamos el concepto de fuerza neta y nos enfocamos en la interacción individual entre una persona y la soga, gobernada por la Tercera Ley de Newton. Las fuerzas actúan en pares. Si la persona (objeto A) ejerce una fuerza de "acción" de $300\text{ N}$ jalando hacia sí misma la soga (objeto B), la soga debe ejercer una fuerza de "reacción" exactamente igual en magnitud y opuesta en dirección sobre las manos de la persona. La soga transmite la tensión a lo largo de su estructura molecular, provocando que tire de cada oponente con la misma intensidad con la que es jalada.

\subsubsection*{Diagrama de Interacciones Concurrentes}
El script a continuación grafica los vectores de fuerza. Muestra claramente la diferencia entre las fuerzas externas aplicadas sobre la soga (cuyo balance es nulo) y las fuerzas reactivas que la soga aplica sobre los extremos externos.

\begin{lstlisting}[language=Python]
import matplotlib.pyplot as plt

fig, ax = plt.subplots(figsize=(8, 4))

ax.plot([-3, 3], [0, 0], color='brown', linewidth=6)

ax.annotate("", xy=(5, 0.5), xytext=(3, 0.5), arrowprops=dict(facecolor='blue', width=3))
ax.text(3, 0.8, "Jalon Persona Der = 300 N", color="blue")
ax.annotate("", xy=(-5, 0.5), xytext=(-3, 0.5), arrowprops=dict(facecolor='red', width=3))
ax.text(-5, 0.8, "Jalon Persona Izq = 300 N", color="red")

ax.annotate("", xy=(1, -0.5), xytext=(3, -0.5), arrowprops=dict(facecolor='cyan', width=3))
ax.text(1.2, -0.9, "Reaccion a Der = 300 N", color="cyan")
ax.annotate("", xy=(-1, -0.5), xytext=(-3, -0.5), arrowprops=dict(facecolor='magenta', width=3))
ax.text(-4.8, -0.9, "Reaccion a Izq = 300 N", color="magenta")

ax.set_xlim(-6, 6)
ax.set_ylim(-2, 2)
plt.title("DCL: Concurso de Tirar la Cuerda")
plt.axis("off")
plt.show()
\end{lstlisting}

\subsubsection*{Conclusión}
La fuerza neta sobre la soga es de $0\text{ N}$. La soga ejerce una fuerza reactiva de $300\text{ N}$ sobre cada persona.

\newpage

\subsection*{Enunciado}
Dos fuerzas actúan sobre un paracaidista que cae en el aire: la fuerza de gravedad y la resistencia del aire. Si la caída es constante, sin ganancia ni pérdida de rapidez, entonces el paracaidista está en equilibrio dinámico. ¿Cómo se comparan las magnitudes de la fuerza gravitacional y de la resistencia del aire?

\subsection*{Solución}

\subsubsection*{Fundamentos del Equilibrio Dinámico y Velocidad Terminal}
Cuando Paul G. Hewitt introduce las Leyes de Newton, hace un énfasis crucial en que la aceleración no es causada por el simple movimiento, sino por un cambio en el estado de movimiento (una fuerza neta desequilibrada). El enunciado estipula que la caída es "constante, sin ganancia ni pérdida de rapidez". En física cinemática, esto significa que la aceleración del paracaidista es exactamente cero ($a = 0$). Este estado particular en la dinámica de fluidos se denomina "velocidad terminal".

\subsubsection*{Igualdad Matemática de Fuerzas Opuestas}
Aplicando la Segunda Ley de Newton ($\Sigma F = m \cdot a$), si la aceleración es cero, el producto de la masa por la aceleración es cero, lo que dicta que la sumatoria de todas las fuerzas actuantes debe ser nula ($\Sigma F = 0$). 
En este escenario unidimensional, la fuerza que tira hacia abajo es el peso del paracaidista (atracción gravitacional) y la fuerza que empuja hacia arriba es la fricción aerodinámica (resistencia del aire). Para que la resta vectorial de estas dos fuerzas colineales y opuestas dé como resultado cero, es un requerimiento matemático absoluto que ambas posean exactamente el mismo valor numérico escalar.

\subsubsection*{Representación Vectorial de la Velocidad Terminal}
Este bloque de código genera el diagrama de cuerpo libre del paracaidista al alcanzar su velocidad límite. Ilustra de forma didáctica que, independientemente de la velocidad de la caída, el reposo vectorial (fuerzas igualadas) es idéntico al que tendría si estuviera colgado de una soga inmóvil.

\begin{lstlisting}[language=Python]
import matplotlib.pyplot as plt

fig, ax = plt.subplots(figsize=(4, 6))

ax.plot(0, 0, 'ko', markersize=20)

ax.annotate("", xy=(0, -6), xytext=(0, -0.5), arrowprops=dict(facecolor='blue', width=4))
ax.text(0.3, -4, "Gravedad (Peso)", color="blue", fontsize=12)

ax.annotate("", xy=(0, 6), xytext=(0, 0.5), arrowprops=dict(facecolor='red', width=4))
ax.text(0.3, 3, "Resistencia del Aire", color="red", fontsize=12)

ax.text(-2.5, 0, "Rapidez Constante\n a = 0 m/s^2", fontsize=11, style='italic')

ax.set_xlim(-3, 3)
ax.set_ylim(-8, 8)
plt.title("Equilibrio Dinamico en Caida Libre")
plt.axis("off")
plt.show()
\end{lstlisting}

\subsubsection*{Conclusión}
Las magnitudes de la fuerza gravitacional y de la resistencia del aire son exactamente iguales.

\newpage

\subsection*{Enunciado}
Un niño aprende en la escuela que la Tierra viaja a más de $100,000\text{ km/h}$ alrededor del Sol y, en un tono alarmado, pregunta por qué no sale disparado. ¿Cuál es tu explicación?

\subsection*{Solución}

\subsubsection*{La Ley de la Inercia y el Marco de Referencia}
La inquietud del niño fue un argumento históricamente utilizado para refutar el modelo heliocéntrico de Copérnico. La explicación moderna, consolidada por Galileo Galilei y estructurada en la Primera Ley de Newton, radica en el concepto de inercia: "Todo objeto continúa en su estado de reposo o de movimiento uniforme en línea recta, a menos que sea obligado a cambiar ese estado por fuerzas aplicadas sobre él". 

\subsubsection*{El Movimiento Compartido}
El niño, los océanos, la atmósfera, los pájaros y todo lo que existe sobre la superficie terrestre están viajando a esos mismos $100,000\text{ km/h}$. Debido a que nosotros y la Tierra compartimos exactamente el mismo estado de movimiento inicial, no necesitamos de una fuerza adicional constante para seguir moviéndonos a esa velocidad; la inercia garantiza que la mantendremos naturalmente. Al igual que cuando saltas dentro de la cabina de un avión que viaja a $900\text{ km/h}$, no te estrellas contra la parte trasera del avión porque tú ya te estás moviendo a $900\text{ km/h}$ junto con él.

\subsubsection*{Ausencia de Aceleración Perceptible}
El ser humano no posee sensores biológicos para detectar la "velocidad absoluta", solo poseemos un sistema vestibular (el oído interno) que detecta la aceleración (los cambios bruscos de velocidad o dirección). Aunque la trayectoria de la Tierra es circular y, por ende, experimenta una aceleración centrípeta hacia el Sol, el radio de la órbita terrestre es tan colosal y la rotación tan suave que el cambio de dirección por segundo es infinitesimalmente pequeño. Esta aceleración es imperceptible para nuestros sentidos y la gravedad de la Tierra es inmensamente más fuerte, manteniéndonos anclados firmemente a su superficie y evitando que salgamos disparados.

\subsubsection*{Conclusión}
No sale disparado debido a la ley de la inercia. Todo lo que está en la Tierra se mueve junto con ella a la misma velocidad constante. Como no experimentamos cambios bruscos de velocidad (aceleración), nuestros sentidos no pueden percibir este movimiento continuo y la gravedad local nos mantiene firmemente unidos al suelo del planeta.

\newpage



\end{document}